% concatenated from hilbert_cubes_BEU_MathComp_submission.tex
\documentclass[10pt]{amsart}
\usepackage{color}
\usepackage{amssymb}
\usepackage{epsfig}
\usepackage{url}
\usepackage{setspace}
\usepackage{algorithm}
\usepackage{algorithmicx}
\usepackage{algpseudocode}
\usepackage{pdflscape}

\theoremstyle{plain}

\newtheorem{thm}{Theorem}[section]
\newtheorem{cor}[thm]{Corollary}
\newtheorem{lem}[thm]{Lemma}
\newtheorem{prop}[thm]{Proposition}
\newtheorem{rem}[thm]{Remark}
\newtheorem{ques}[thm]{Question}
\newtheorem{conj}[thm]{Conjecture}
\newtheorem{exam}[thm]{Example}
\newtheorem{prob}[thm]{Problem}
\def\cal{\mathcal}
\def\bbb{\mathbb}
\def\op{\operatorname}
\renewcommand{\phi}{\varphi}
\newcommand{\R}{\bbb{R}}
\newcommand{\N}{\bbb{N}}
\newcommand{\Z}{\bbb{Z}}
\newcommand{\Q}{\bbb{Q}}
\newcommand{\F}{\bbb{F}}
\newcommand{\ep}{\epsilon}
\newcommand{\tht}{\theta}
\newcommand{\cir}{\bbb{S}^1}
\newcommand{\C}{\bbb{C}}
\newcommand{\eps}{\varepsilon}
\newcommand{\aand}{\text{\rm\ \ and\ \ }}
\newcommand{\Xk}{X^{*k}}
\newcommand{\ied}{I^{\textrm{ED}}}
%\setlength{\oddsidemargin}{0.5cm}  %Print setup: top: -1cm
%\setlength{\evensidemargin}{0.5cm} %Left margin: 2.5cm
%\setlength{\topmargin}{-2.0cm} \setlength{\textheight}{25cm}
%\setlength{\textwidth}{15cm}

\begin{document}
\title[Hilbert cubes of dimension 3 in the set of squares]{There are infinitely many Hilbert cubes of dimension 3 in the set of squares}

\author{Andrew Bremner, Christian Elsholtz and Maciej Ulas}

\keywords{Hilbert cube, sumsets in the squares, rational curve, elliptic curve, parameterization, Bombieri-Lang conjecture}
\subjclass[2000]{11D41, 11G05, 11B75, 14H52}
\thanks{
Research of C.E. was supported by a joint FWF-ANR project ArithRand (I 4945-N and ANR-20-CE91-0006).\\
Research of M.U. was supported by a grant of the National Science Centre (NCN), Poland, no. UMO-2019/34/E/ST1/00094}

%\today


\begin{abstract}
A Hilbert cube of dimension $d$ is the set of integers 
\[
H(a_{0}; a_{1}, \ldots, a_{d})=a_{0}+\{0, a_{1}\}+\cdots+\{0, a_{d}\}=\left\{a_{0}+\sum_{i=1}^{d}\varepsilon_{i}a_{i}:\;\varepsilon_{i}\in\{0,1\}\right\}.
\]
Brown, Erd\H{o}s and Freedman asked whether the maximal dimension of a Hilbert cube in the set $\cal{S}=\{n^2:\;n\in\mathbb{N}\}$ of integer squares is absolutely bounded or not. Dietmann and Elsholtz proved that if $H(a_{0}; a_{1}, \ldots, a_{d})\subset \cal{S}\cap [0, N]$, then $d\leq 7 \log\log N$ for all sufficiently large values of $N$. Here we prove that there exist at least $\gg N^{1/8}$ Hilbert cubes 
$H(a_{0}; a_{1}, a_{2}, a_{3})$ with $a_{0}, a_{1}, a_{2}, a_{3}\in [0,N]$ in the set of squares.
Moreover, we prove that for each $i, j\in\{0, 1, 2, 3\}$ with $i<j$, the set
$$
\left\{\frac{a_{i}}{a_{j}}:\;H(a_{0}; a_{1}, a_{2}, a_{3})\subset S\right\}
$$
is dense in the set of positive real numbers (in the Euclidean topology).
\end{abstract}
\maketitle
\section{Introduction}\label{sec1}



A Hilbert cube 
\[H(a_0;a_1, \ldots, a_d)=a_0+ \{0,a_1\}+ \cdots + \{0,a_d\}
=\left\{a_{0}+\sum_{i=1}^{d}\varepsilon_{i}a_{i}:\;\varepsilon_{i}\in\{0,1\}\right\}
\]
is an iterated sumset. These cubes were introduced by Hilbert \cite{Hilbert:1892} in connection with an irreducibility question on polynomials.
These cubes occur quite naturally for example in Szemer\'{e}di's proof of
Szemer\'{e}di's theorem \cite{Szemeredi:1969}, for arithmetic progressions of length $4$ and in Gowers's norm.
They are also a well studied object in combinatorics, see e.g.~\cite{Conlon-Fox-Sudakov:2014, Hegyvari:1999, Solymosi:2007}.

For a given set $S$ of positive integers it is natural to ask if infinite Hilbert cubes can exist in $S$,
and, if not, about the maximal dimension $d$ of such a cube in $S$. It is easy to see that an infinite Hilbert cube does not exist in the squares, as for $i \geq 1$ each $a_i$ would be the difference between squares infinitely often, 
(say $a_i=(a_0+a_i+a_j)-(a_0+a_j)$). But as
differences between consecutive squares are increasing, a fixed difference can be the difference between squares only finitely often.

Let us assume that all $a_1, \ldots , a_d$ are non-zero.
Brown, Erd\H{o}s and Freedman \cite{BrownandErdosandFreedman:1990}
asked whether the maximal dimension of a Hilbert cube in the set $\cal{S}=\{n^2:\;n\in\mathbb{N}\}$ of integer squares is \emph{absolutely bounded} or not. 
Hegyvari and S\'{a}rkozy \cite{HegyvariandSarkozy:1999}
 proved for the set of squares in a finite interval $[1,N]$
that the maximal value of $d$ is bounded by $d=O((\log N)^{1/3})$.
This was improved by Dietmann and Elsholtz \cite{Dietmann-Elsholtz:2012, Dietmann-Elsholtz:2015} first to $d=O((\log \log N)^2)$, and later to $O(\log \log N)$.

Cilleruelo and Granville
\cite{CillerueloandGranville:2007} observed that
the Bombieri-Lang conjecture implies that $d$ is absolutely bounded.
In fact, $x^2+a_1, x^2+a_2$ and $x^2+a_1+a_2$ must be squares for the many distinct values of $x=a_0 + \{0,a_3\} + \cdots + \{0,a_d\}$, while the Bombieri-Lang conjecture implies that the curve
$y^2=(x^2+a_1)(x^2+a_2)(x^2+a_1+a_2)$ only has an absolutely bounded number of integral solutions $(x,y)$,
(for details see \cite{CaporasoHarrisandMazur, ShkredovandSolymosi}).

Some related problems are as follows:
A well known open problem by Rudin asks how many squares 
can be in an arithmetic progression $an+b, 1\leq n \leq N$. Rudin conjectured that the progression $24n+1$ has the largest number of squares. For large modulus this is a very difficult question.
Bombieri, Granville and Pintz \cite{BombieriGranvilleandPintz},  and later Bombieri and Zannier \cite{BombieriandZannier} proved that the  maximal number of squares in an arithmetic progression of length $N$ is bounded by $O(N^{2/3}(\log N)^A))$ and $O((\log N)^{3/5}(\log N)^{A'})$, respectively, where $A,A'$ are constants.

One may ask about sumsets $A+B$ in the set of squares. A Hilbert cube of high dimension 
can be naturally decomposed into two such summands, in many distinct ways.
Erd\H{o}s and Moser asked whether
there are arbitrarily large sets
such that $a_i+a_j$ are always squares, where $a_i, a_j \in A$ are distinct. 
Rivat, S\'{a}rk\'{o}zy and Stewart \cite{RivatSarkozyandStewart} and 
Gyarmati \cite{Gyarmati:2001} studied the size of sumsets in the squares $\cal{A}+\cal{B}\subset \cal{S}\cap [1,N]$.
For example, the following bound holds (see \cite{Gyarmati:2001}):  
$\min (|\cal{A}|, \cal{B}|)\leq 8 \log N$.
Related results are by the second author with Wurzinger \cite{Elsholtz-Wurzinger}, for example:
If $|\cal A|\geq \log \log N$, then $|\cal B| =O(\log N \log \log N)$.
The second author with Ruzsa and Wurzinger \cite{Elsholtz-Ruzsa-Wurzinger} proved:
If $\cal{A}+\cal{B}+\cal{C}\subset \cal{S}\cap [1,N]$, then
$\min (|\cal{A}|, \cal{B}|, |\cal{C}|)=O(( \log N)^{4/5})$.


When $|A|=3$ the situation can be well described by elliptic curves, see \cite{DujellaandElsholtz:2013}.
In fact, for every sufficiently large $N$ there exist sets $A$ such that there exists a set $B$ of size 
$\gg (\log N)^{15/17}$ such that all elements of $A+B\subset [1,N]$ are in the squares. Further, it is also possible to have all elements $A,B,A+B$ being squares, with the estimate
$\gg (\log N)^{9/11}$.

When $|A|=4$ or larger,
then conditional on the Bombieri-Lang conjecture (or the uniformity conjecture) there is an {\emph{absolute bound}} on $|B|$.
(See section 4.3. of \cite{AlonandAngelandBenjaminiandLubetzky:2012}.)

For comparison it is known that the set of squares does not contain any arithmetic progression of length $4$.


As it turns out, one can find some small Hilbert cubes by hand, and many more by a computer.
Two  easy examples are
\[1+\{0,15\}+\{0,48\}=\{1^2,4^2,7^2,8^2\}\]
and
\[1 + \{0, 528\} + \{0, 840\}+ \{0, 840\} = \{1^2, 23^2, 29^2, 29^2, 37^2, 37^2, 41^2, 47^2\}.\]
Note that in the last example $29^2$ and $37^2$ occur in two different ways as the set of squares contains \emph{two} arithmetic progressions of length $3$ with gap 840.
According to our computations the smallest Hilbert cube of dimension 3 with all entries positive and distinct is 
$$
10^2+\{0, 2400\}+\{0, 4389\}+\{0, 8736\}=\{10^2, 50^2, 67^2, 83^2, 94^2, 106^2, 115^2, 125^2\}.
$$
Here, by smallest, we mean the Hilbert cube $H(a_{0}; a_{1}, a_{2}, a_{3})$ with the smallest possible sum $a_{0}+a_{1}+a_{2}+a_{3}$.

We will concentrate on the case that the $a_i$ are positive.
Otherwise one could have closely related Hilbert cube representations such as
\begin{alignat*}{3}
2209 + \{0, -528\} + \{0, -840\}+ \{0, -840\} &=&
841 + \{0, 528\} + \{0, -840\}+ \{0, 840\}\\ 
&=&
529 + \{0, -528\} + \{0, 840\}+ \{0, 840\}\\
&=&\{1^2, 23^2, 29^2, 29^2, 37^2, 37^2, 41^2, 47^2\}.
\end{alignat*}

It might seem that a $d$-dimensional Hilbert cube with many overlapping expressions, for example $a_0+a_1+a_2+a_3=a_0+a_4$ might be easier to find as there are less square-conditions. A natural question is: how many distinct elements must a $d$-dimensional Hilbert cube in the squares necessarily have?
Dietmann and the second author \cite[Lemma 1.4]{Dietmann-Elsholtz:2015} gave an exponential lower bound:
$|\cal{H}(a_0;a_1, \ldots, a_d)|\geq \lceil 2 (\frac{4}{3})^{d-1}\rceil$.
That bound was intended for the asymptotic exponential growth.
For small $d$ it can be adapted as follows:
\begin{lem}
Let $k \geq 3$ be a positive integer, and let $\cal{S}$ denote a set of integers without an arithmetic progression of length $k$. 
Moreover, let $H(a_0; a_1, \ldots , a_d) \subset \cal{S}$. Let $d=d_1+d_2$. Then
\[|H| \geq \left\lceil C_{d_1}  (\frac{k}{k-1})^{d_2}\right\rceil, \]
where $C_{d_1}$ is the minimal size of a Hilbert cube of dimension $d_1$ in ${\cal S}$.
\end{lem}
It is well known that for sets without progressions of length $k= 3$ one even has that $|H| = 2^d$. For the set of squares we have $k=4$.

Let $d=1$, then the estimate gives the correct value $C_1=2$. When $d=2$ this gives the correct value $C_2=\lceil\frac{8}{3}\rceil=3$. 
When $d=3$ the lower bound gives $4$, while $6$ is the actual minimum. 
%When $d=4$ this gives the value $6$.
%We would expect the minimum to be either $9$, or $12$.

%\Comment{I am not so certain about the guesses 9 and 12 above. Maybe we better remove this?, I believe the minimum 6 for $d=3$ can be explained.}


Here are several quite natural open problems.
\begin{ques}\label{firstques}
 \begin{enumerate}
\item
Does a Hilbert cube with $d=4$ exist in the set of squares? Do infinitely many such cubes exist? 
\item Does a Hilbert cube
with $a_0=0$ and $d=3$ exist?
(This is the open problem of a perfect Euler brick (or cuboid).)
\item
 Can the value of $d$ increase (without any upper bound)?
\item
If the last question has a negative answer:
What is the largest value of $d$ which occurs at least once,
and what is the largest value of $d$ which occurs for infinitely many Hilbert cubes in $\cal{S}$?
\item
Similarly, what is the largest possible value of $d$, when all $a_i$ are distinct?
\end{enumerate}
\end{ques}

For all we know no Hilbert cube has been found in the set of squares of dimension $d \geq 4$.
But even for dimensions $d=2$ and $d=3$ there is no systematic study.
For dimension $d=2$ it is not difficult to write down solutions:
A Hilbert cube $H(a_0;a_1,a_2)=\{a_0, a_0+a_1, a_0+a_2, a_0+a_1+a_2\}$ in the set of squares can be described by
$a_0=x^2, a_1=y^2-x^2, a_2=z^2-x^2$, if and only if
$a_0+a_1+a_2=y^2+z^2-x^2$ is a square.
In other words, every integer with at least two distinct representations as sums of two squares, $y^2+z^2=t^2+x^2$, generates a Hilbert cube of dimension $2$ in the squares, and every Hilbert 
cube of dimension $2$ corresponds to a solution of 
$y^2+z^2=t^2+x^2$.
Moreover, if $a_1=a_2$, then $y=z$ and this corresponds to $2y^2=x^2+t^2$, i.e.~three distinct squares in an arithmetic progressions. (Note that the case $x=t$ would yield $x=y=t$, and so $a_1=x^2, a_1=y^2-x^2=0, a_2=z^2-x^2=0$).

A Hilbert cube is called reduced if $\gcd(a_0,a_1,a_2,a_3)=1$.
In this manuscript we show that infinitely many reduced Hilbert cubes
of dimension $d=3$ exist, giving an explicit parametric construction. This leads to a good quantitative lower bound for the number of such cubes in $[0,N]$.

%\Comment{Do we dare to make a conjecture about an almost best possible lower bound? Any heuristic possible?}

Looking at the table of examples many further questions arise with regard to the values $a_i$ or patterns:
\begin{ques}\label{genques}
\begin{enumerate}
    \item 
    Can there be infinitely many Hilbert cubes with a fixed value of $a_0$? Otherwise, can one give an effective upper bound on the number and the size of elements $a_1,a_2,a_3$ in terms of $a_0$?
     \item Are there infinitely many reduced Hilbert cubes of dimension $d \geq 3$ where
     $a_i=a_{i+1}$ holds? (The cube will then contain an arithmetic progression of length 3, and many shifted copies.)
\item Does a Hilbert cube of dimension $4$ with $a_1=a_2>0$ and $ a_3=a_4>0$ exist? Note that if so, then we may construct a $3 \times 3$ magic square of squares, with common sum $3(a_0+a_1+a_3)$ of rows, columns and diagonals:
\begin{center}
        $\Biggl($
\begin{tabular}{ccc}
        $a_0+2 a_1+a_3$ & $a_0$ & $a_0+a_1+2 a_3$ \\
        $a_0+2 a_3$ & $a_0+a_1+a_3$ & $a_0+2 a_1$ \\
        $a_0+a_1$ & $a_0+2 a_1+2 a_3$ & $a_0+a_3$
\end{tabular}
        $\Biggr)$
\end{center}
This special case is therefore a well-known open question: see \cite{Robertson}, \cite{Bremner:1999}, \cite{Bremner:2001}. The magic square (see \cite{Bremner:1999})
\begin{center}
        $\Biggl($
\begin{tabular}{ccc}
        $565^2$ & $23^2$ & $222121$ \\
        $289^2$ & $425^2$ & $527^2$ \\
        $373^2$ & $360721$ & $205^2$
\end{tabular}
        $\Biggr)$
\end{center}
leads to $(a_0,a_1,a_2,a_3,a_4)=(23^2, \; 138600, \; 138600, \; 41496, \; 41496)$ with thirteen of the sixteen required sums for a Hilbert cube of dimension $4$ being square.
     \item
     Given $a_0$, or $a_0,a_1$ or $a_0,a_1,a_2$ what are good search bounds 
     on other elements $a_i$ and how can one algorithmically find cube extensions?
     \end{enumerate}
     \end{ques}

We address some of these questions in this paper and give the following 
results.

\begin{thm}{\label{thm:finite}}
\begin{enumerate}
\item  For any fixed square $a_0$ there are infinitely many Hilbert cubes in the set of squares with this $a_0$ and $d=2$.
\item For fixed $a_1$ there are only finitely many Hilbert cubes of 
dimension $d \geq 2$ in the set of squares. The bound is effective:
\begin{enumerate}
\item
For each $a_i$ the upper bound $a_i\leq \frac{1}{4}(a_1-1)^2-a_0$ holds.
\item
For each $a_i$ there are at most $\tau(a_i)$ many choices, where $\tau(n)$ is the number of positive divisors of $n$.
\end{enumerate}
\end{enumerate}
\end{thm}

Let us introduce the set
$$
\cal{H}:=\{(a_{0}, a_{1},a_{2},a_{3})\in\N\times \Z^{3}:\;H(a_{0}; a_{1}, a_{2}, a_{3})\subset \cal{S}\}
$$
containing all Hilbert cubes of dimension three, sitting in the set of squares.
\begin{thm}{\label{thm:H3}}
The number $H_3(N)$ of reduced Hilbert cubes in $\cal{H}$
of dimension $d=3$ and $0\leq a_i\leq N$, for $i=0, \ldots, d$ satisfies
$$
H_3(N)\gg N^{1/8}.
$$
\end{thm}


\begin{thm}{\label{thm:idenical base}}
There exist infinitely many reduced Hilbert cubes in ${\cal H}$  of dimension $d=3$, with
$a_1=a_2<a_3$,    see equation (\ref{a1a2solution}). 
Moreover, the number of such Hilbert cubes with $a_i \leq N$ 
is at least $cN^{1/10}$, for some positive constant $c$.
The same statement holds for $a_1<a_2=a_3$.

\end{thm}


\begin{thm}{\label{thm:H_d-upper-bound}}
The number $H_2(N)$ of Hilbert cubes (in the set of squares) 
of dimension $2$ and $0\leq a_i\leq N$, for $i=0,1,2$ satisfies
\[N\log N \ll H_2(N) \ll N\log N.\]
The number $H_d(N)$ of Hilbert cubes of dimension $d$ and $0\leq a_i\leq N$, for $i=0, \ldots, d$ satisfies, for all $\varepsilon_d>0$,
\[H_d(N) \ll_d N^{1+\varepsilon_d}.\]
\end{thm}

From the parametric family of Hilbert cubes 
in dimension $3$ one can deduce the following:
\begin{cor}{\label{cor-fractions}}
Let $\cal{H}_{+}=\cal{H}\cap \N_{+}^{4}$. For each $i, j\in\{0, 1, 2, 3\}$ with $i<j$, the set
$$
H_{i,j}:=\left\{\frac{a_{i}}{a_{j}}:\;(a_{0}, a_{1}, a_{2}, a_{3})\in\cal{H}_{+}\right\}\subset \Q_{+}
$$
is dense in the Euclidean topology in the set $\R_{+}$.
\end{cor}



\begin{proof}[Proof of Theorem \ref{thm:H_d-upper-bound}]
If $a_{0}=p^2, a_{0}+a_{1}=q^2$, then there are at most $\sqrt{N}$ choices for each of $a_0\leq N$ and $a_1\leq N$.
Having fixed $a_1=q^2-p^2$, the number of choices of $a_2, a_3, \ldots , a_d$ is bounded by the number of divisors $\tau(a_1)$ each, which easily gives the upper bound: for every $\epsilon>0$ 
$H_d(N)\leq N\exp(c_d \frac{\log N}{\log \log N})=O(N^{1+\varepsilon})$.
It should be possible to find a stronger upper bound.
\end{proof}

The question about Hilbert cubes in the set of quadratic residues has been studied by
\cite{Alsetri-Shao, Dietmann-Elsholtz:2012,
Dietmann-Elsholtz-Shparlinski,
HegyvariandSarkozy:1999, Schoen}.


A closely related question is about sumsets $A+B$ contained in the set of squares. There is some information 
in \cite{AlonandAngelandBenjaminiandLubetzky:2012, DujellaandElsholtz:2013}.



%Two Hilbert cubes $H_1, H_2$ are affinely equivalent, if one can write %$H_2=b_0+qH_1=(b_0+qa_0)+\{0,qa_1\}+\cdots +\{0,q a_d\}$.


\section{Notation}
Let us introduce the following notation.

\begin{equation}\label{Hsystem}
\begin{aligned}
p^2&=a_{0},\\
q^2&=a_{0}+a_{1},\\
r^2&=a_{0}+a_{2},\\
s^2&=a_{0}+a_{1}+a_{2},\\
P^2&=a_{0}+a_{3},\\
Q^2&=a_{0}+a_{1}+a_{3},\\
R^2&=a_{0}+a_{2}+a_{3},\\
S^2&=a_{0}+a_{1}+a_{2}+a_{3},
\end{aligned}
\end{equation}
for some integers $p, q,r,s, P,Q,R, S$.

\section{Small solutions}\label{sec2}
As usual, facing a Diophantine problem it is natural to look for small integer solutions using  a computational approach. However, before we present the results of our initial Let us note that there are some natural maps which act on the set $\cal{H}$. Indeed, each of the following maps  
\begin{align*}
\phi_{1}&:\;\cal{H}\ni (a_{0}, a_{1}, a_{2}, a_{3})\mapsto (a_{0}+a_{1}, -a_{1}, a_{2}, a_{3})\in \cal{H},\\
\phi_{2}&:\;\cal{H}\ni (a_{0}, a_{1}, a_{2}, a_{3})\mapsto (a_{0}+a_{2}+a_{3}, a_{1}, -a_{3}, -a_{2})\in \cal{H},\\
\phi_{\sigma}&:\;\cal{H}\ni (a_{0}, a_{1}, a_{2}, a_{3})\mapsto (a_{0}, a_{\sigma(1)}, a_{\sigma(2)}, a_{\sigma(3)})\in \cal{H},
\end{align*}
where $\sigma$ is an element of the permutation group $\Sigma_{3}$ of three elements, act on $\cal{H}$. We have $\phi_{i}^{(2)}=\op{id}_{\cal{H}}$ for $i=1, 2$ and for each $\sigma\in \Sigma_{3}$ we have $\phi_{\sigma}^{(6)}=\op{id}_{\cal{H}}$. We thus see that there is a group of order $2\times 2\times 6=24$ generated by the above maps which acts on $\cal{H}$. Moreover, for each positive integer $m$ we have an additional map
$$
\psi_{m}:\;\cal{H}\ni (a_{0}, a_{1}, a_{2}, a_{3})\mapsto  (m^2a_{0}, m^2a_{1}, m^2a_{2}, m^2a_{3})\in\cal{H}.
$$
As a consequence, after applications of suitable compositions of maps, without loss of generality we can assume that $0< a_{1}\leq a_{2}\leq a_{3}$ and $\gcd(a_{0}, a_{1}, a_{2}, a_{3})=1$ and obtain a reduced tuple.

First of all, let us start with the following simple
\begin{lem}
Let $m\in\N_{\geq 2}$ be fixed. Then, for any given pair of positive integers $a_{0}, a_{1}$, with $a_{0}<a_{1}$ there are only finitely many Hilbert cubes of dimension $m$ in the set of squares.
\end{lem}
\begin{proof}
We start with the case $m=2$. We have that $H(a_{0}; a_{1}, a_{2})\subset \cal{S}$ if and only if
$$
a_{0}=p^2,\quad  a_{0}+a_{1}=q^{2}, \quad  a_{0}+a_{2}=r^2, \quad a_{0}+a_{1}+a_{2}=s^2
$$
for some $p, q, r, s$. Because $a_{0}, a_{1}$ are fixed, $p, q$ are also fixed. Subtracting the third equation from the fourth we get that $a_{1}=s^2-r^2=(s-r)(s+r)$. In other words, if $d$ is a divisor of $a_{1}$ then the corresponding values of $r, s$ are of the form
$$
r=\frac{1}{2}\left(\frac{a_{1}}{d}-d\right),\quad s=\frac{1}{2}\left(\frac{a_{1}}{d}+d\right).
$$
Having possible $r, s$ we take $a_{2}=r^2-a_{0}$. 
Because the number of possible values of $r$ is bounded by $\tau(a_{1})$, where $\tau(n)$ is the number of divisors of $n$, the number of possible values of 
$a_{2}$ is also bounded by $\tau(a_{1})$. This proves the lemma for $m=2$. However, by iterating this argument, we see that for any given $a_{0}, a_{1}$ and $m$, the number of Hilbert cubes $H(a_{0}; a_{1}, a_{2},\ldots, a_{m})$ contained in the set of squares is finite.
\end{proof}

Although very simple, the method from the proof of the above lemma can be used to write an efficient algorithm which allows to find all Hilbert cubes $H(a_{0}; a_{1}, a_{2}, a_{3})\subset \cal{S}$ with $a_{0}, a_{1}\leq N$ for given $N$ in time $O_{\varepsilon}(N^{1+\varepsilon})$
The algorithm runs as follows. First of all, for a given $a_{1}$ we compute the set 
$$
D_{small}(a_{1})=\{d\in\N:\;d|a_{1}\;\mbox{and}\; d^2<a_{1}\}=\{d_{1}, d_{2}, \ldots, d_{m}\},
$$
in natural order,
where $m=m(a_{1})=\#D_{small}(a_{1})$. Next, for given $d_{i}, d_{j}\in D_{small}(a_{1}), i<j$, we check whether 
\begin{align*}
&c_{1}(i,j)=\frac{1}{4}\left(\frac{a_{1}}{d_{i}}-d_{i}\right)^2
+\frac{1}{4}\left(\frac{a_{1}}{d_{j}}-d_{j}\right)^2-a_{0}=\square,\\
&c_{2}(i,j)=a_{1}+\frac{1}{4}\left(\frac{a_{1}}{d_{i}}-d_{i}\right)^2
+\frac{1}{4}\left(\frac{a_{1}}{d_{j}}-d_{j}\right)^2-a_{0}=\square.
\end{align*}

If these conditions are satisfied, then we put 
$$
a_{2, i}=\frac{1}{4}\left(\frac{a_{1}}{d_{i}}-d_{i}\right)^2-a_{0},\quad a_{3, j}=\frac{1}{4}\left(\frac{a_{1}}{d_{j}}-d_{j}\right)^2-a_{0}.
$$
and, provided that $H(a_{0}; a_{1})\subset \cal{S}$, we get a Hilbert cube $H(a_{0}; a_{1}, a_{2,i}, a_{3, j})\subset \cal{S}$. 

The pseudo-code of our algorithm has the following form. We implemented it in Pari GP \cite{PARI2}.

\begin{algorithm}[H]
\caption{Finding Hilbert cubes $H(a_{0};a_{1},a_{2},a_{3}) \subset \mathcal{S}$ with $a_{0}<a_{1}\leq N$}
\begin{algorithmic}[1]
\Procedure{HilbertCubes}{$N$}
  \For{$p = 1$ \textbf{to} $\sqrt{N}$; $a_0=p^2$}
        \For{$a_{1} = a_{0}+1$ \textbf{to} $N$}
      \If{$a_{0}+a_{1}$ is not a square} \State \textbf{continue} \EndIf
      \State $D \gets \{ d \in \mathbb{N} : d \mid a_{1}, \; d^{2}<a_{1}\}$
      \ForAll{pairs $(d_i,d_j)$ with $i<j$ from $D$}
        \State $r_i \gets \tfrac{1}{2}\left(\tfrac{a_{1}}{d_i} - d_i\right)$
        \State $r_j \gets \tfrac{1}{2}\left(\tfrac{a_{1}}{d_j} - d_j\right)$
        \If{$r_i$ or $r_j$ not integer} \State \textbf{continue} \EndIf
        \State $a_{2} \gets r_i^{2} - a_{0}$, \quad $a_{3} \gets r_j^{2} - a_{0}$
        \If{$a_{2} \leq 0$ or $a_{3} \leq 0$} \State \textbf{continue} \EndIf
        \If{$a_{0}+a_{2},\; a_{0}+a_{3},\; a_{0}+a_{2}+a_{3},\;
              a_{0}+a_{1}+a_{2},\; a_{0}+a_{1}+a_{3},\; a_{0}+a_{1}+a_{2}+a_{3}$ are all squares \textbf{and} $\gcd(a_0, a_1, a_2, a_3)=1$}
          \State \textbf{output} $[a_{0},\op{sort}(a_{1},a_{2},a_{3})]$
        \EndIf
      \EndFor
    \EndFor
  \EndFor
\EndProcedure
\end{algorithmic}
\end{algorithm}


\begin{exam}
{\rm To see the described algorithm in action let us consider the following examples. 
Take $a_{0}=1, a_{1}=8099=90^2-1$, i.e., $H(a_{0}; a_{1})\subset \cal{S}$. We have 
$$
D_{small}(a_{1})=\{1, 7, 13, 89\}=\{d_{1}, d_{2}, d_{3}, d_{4}\}.
$$
The corresponding values of $c_{1}(i, j), c_{2}(i, j)$ are presented in tables below. 
\begin{table}[htbp]
\centering
\begin{tabular}{|l|llll|}
\hline
 $i\backslash j$ & 1       & 2        & 3        & 4\\
 \hline
   1           &32788801 & 16725025 & 16487425 & $4049^2$ \\
    2          &&661249 & 423649 & $575^2$  \\
     3         &&&186049 & $305^2$  \\
      4        &&&&1   \\
      \hline
\end{tabular}
\caption{Values of $c_{1}(i, j), 1\leq i<j\leq 4$ for $a_{0}=1, a_{1}=8099$.}
\end{table}

\begin{table}[htbp]
\centering
\begin{tabular}{|l|llll|}
\hline
 $i\backslash j$ & 1       & 2        & 3        & 4\\
 \hline
1 &32796897 & 16733121 & 16495521 & 16402497 \\
2 &&669345 & 431745 & 338721  \\
3 &&&194145 & 101121  \\
4 &&&&8097 \\
     \hline
\end{tabular}
\caption{Values of $c_{2}(i, j), 1\leq i<j\leq 4$ for $a_{0}=1, a_{1}=8099$}
\end{table}

We note that $c_{1}(i, j)$ is a square if and only if $j=4$. However, the value $c_{2}(i, 4)$ is not a square for any $i=1, 2, 3,4 $. In consequence, there cannot be any Hilbert cube of dimension $d\geq 2$ with $a_{0}=1, a_{1}=8099$.

To see other example let us take $a_{0}=4$ and $a_{1}=3360=58^2-4$. We then have
$$
D_{small}(a_{1})=\{1, 2, 3, 4, 5, 6, 7, 8, 10, 12, 14, 15, 16, 20, 21, 24, 28, 30, 32, 35, 40, 42, 48, 56\}.
$$
Performing all necessary calculations we find that 
$$
c_{1}(13,16)=113^2,\quad c_{2}(13,16)=127^2
$$ 
(related to the divisors $d_{13}=16, d_{16}=24$ of $a_{1}$) and we get the corresponding Hilbert cube of dimension 3 is $\{4,3360,9405,3360\}$. Interestingly, there is another solution given by 
$$
c_{1}(16,17)=74^2,\quad c_{2}(16,17)=94^2
$$ 
(related to the divisors $d_{16}=24, d_{17}=28$ of $a_{1}$) and we get the corresponding Hilbert cube of dimension 3 is $\{4, 2112, 3360,3360\}$.}
\end{exam}

The above result allows us to find all reduced Hilbert cubes of dimension 3 with $a_{0}, a_{1}<10^8$. To do that we used the described algorithm with $a_{0}=p^2, a_{1}=q^2-p^2$ with $p\leq 10^4,  p<q\leq \sqrt{p^2+10^{8}}$. Under these assumptions we found 4644 reduced Hilbert cubes of dimension three sitting in the set of squares (the computations took about one day). None of these cubes can be extended to a Hilbert cube of dimension 4. In the table below we present all three dimensional Hilbert cubes $H(a_{0}; a_{1}, a_{2}, a_{3})$ with $\op{max}\{a_{0}, a_{1}, a_{2}, a_{3}\}\leq 2^{18}.$

In this context it is natural to consider the quantity
$$
C_{3}(N)=\#\{(a_{0}, a_{1})\in [0,N]^2:\;a_{0}<a_{1}, H(a_{0}; a_{1}, a_{2}, a_{3})\subset\cal{S}\;\mbox{for some}\;a_{2}, a_{3}\in\N\}.
$$
In the table below we present the values of $H_{3}(m)$ and $C_{3}(m)$ of the number of reduced Hilbert cubes $H(a_{0}; a_{1}, a_{2}, a_{3})$ with $\op{max}\{a_{0}, a_{1}, a_{2}, a_{3}\}\leq m$ and  $\op{max}\{a_{0}, a_{1}\}\leq m$ respectively, for $m=2^{n}$ and $n=15,\ldots, 26 $.
\begin{table}[htbp]
\centering
\begin{tabular}{|c|c|c|c|c|c|c|c|c|c|c|c|c|}
\hline
 $n$         & 15 & 16 & 17 & 18 & 19 & 20 & 21 & 22 & 23 & 24 & 25 & 26 \\
 \hline 
$H_{3}(2^{n})$  & 8 & 13 & 28 & 39 & 65 & 99 & 147 & 239 & 363 & 529 & 792 & 1173 \\
 \hline
 $C_{3}(2^{n})$   &  27 & 51 & 91 & 136 & 228 & 349 & 541 & 852 & 1278 & 1851 & 2730 & 3887\\
 \hline
\end{tabular}
\caption{Values of $H_{3}(2^n)$ and $C_{3}(2^n)$ for $n\in\{15,\ldots, 26\}$.}
\end{table}
%\begin{figure}[h!]
%    \centering
%    \includegraphics[width=0.9\linewidth]{Pic_1.jpg}
%    \caption{Plot of the values $H_{3}(n)$ for $n\leq 10^7$.}
%    \label{fig:Pic2}
%\end{figure}
\begin{figure}[!ht]
%\begin{figure}[h!]
    \centering
    \includegraphics[width=0.9\linewidth]{Pic_3.jpg}
    \caption{Plot of the functions $C_{3}(n)$ (blue) and $H_{3}(n)$ (yellow) for $n\leq 10^6$.}
    \label{fig:Pic3}
\end{figure}


%Taking into accounts all these observations we computed all %Hilbert cubes $H(a_{0}; a_{1}, a_{2}, a_{3})$ of dimension 3 %with $\op{max}\{a_{0}, a_{1}, a_{2}, a_{3}\}\leq 10^6$. To do %this let us come back to the system (\ref{Hsystem}) and solve %that the (sub)system consisting the first, the second, the third %and the fifth equation with respect to $a_{0}, a_{1}, a_{2}, %a_{3}$. We get
%$$
%a_{0}=p^2, \quad a_{1}=q^2-p^2, \quad a_{2}=r^2-p^2, \quad %a_{3}=P^2-p^2
%$$
%and we are left with the following system:
%\begin{equation*}
%\begin{cases}
%\begin{array}{lll}
% s^2&=&q^2+r^2-p^2 \\
% Q^2&=&q^2+P^2-p^2 \\
% R^2&=&r^2+P^2-p^2  \\
% S^2&=&q^2+r^2+P^2-2p^2  \\
%\end{array}
%\end{cases}
%\end{equation*}
%Looping with $p\leq 10^3, p<q\leq r\leq P\leq \sqrt{10^6+p^2}$ %we covered the range we were interested in. The results of our %computations are presented in the table below.
%{\small 
%\begin{equation*}
%\begin{array}{|llll||llll|}
%\hline
%       a_{0} & a_{1} & a_{2} & a_{3} &     a_{0} & a_{1} & a_{2} & a_{3}\\
%\hline
% 1^2 & 528 & 840 & 840 &                       103^2 & 472416 & 839475 & 897600 \\
%  2^2 & 3360 & 3360 & 9405 &                     107^2 & 18480 & 28152 & 127680 \\
%  2^2 & 4485 & 7392 & 20160 &                    107^2 & 57195 & 364320 & 446880 \\
%  2^2 & 17952 & 26565 & 131040 &                 107^2 & 92235 & 446880 & 646272 \\
%  2^2 & 100485 & 209760 & 376992 &               107^2 & 102120 & 247632 & 414960 \\
%  5^2 & 32736 & 46200 & 54264 &                 109^2 & 19803 & 36960 & 163680 \\
%  5^2 & 78936 & 218064 & 319200 &               120^2 & 106704 & 426496 & 482625 \\
%  7^2 & 5280 & 5280 & 9555 &                    123^2 & 176715 & 235872 & 978880 \\
%  7^2 & 40755 & 108192 & 252960 &               125^2 & 22011 & 84864 & 117600 \\
%  7^2 & 127400 & 180576 & 180576 &              127^2 & 33600 & 46371 & 280896 \\
%  10^2 & 2400 & 4389 & 8736 &                   129^2 & 17955 & 40480 & 157248 \\
%  10^2 & 24864 & 184800 & 591261 &              129^2 & 177840 & 366520 & 506088 \\
%  13^2 & 10440 & 15960 & 62832 &                130^2 & 94656 & 181125 & 421344 \\
%  16^2 & 175305 & 802560 & 929040 &             131^2 & 132608 & 270135 & 720720 \\
%  17^2 & 45936 & 134400 & 249711 &              134^2 & 302400 & 544544 & 730269 \\
%  23^2 & 41496 & 138600 & 138600 &              139^2 & 115368 & 182280 & 318240 \\
%  31^2 & 10920 & 10920 & 26928 &                141^2 & 37240 & 201960 & 201960 \\
%  31^2 & 38640 & 57120 & 99528 &                145^2 & 180576 & 519200 & 943299 \\
%  31^2 & 184800 & 503139 & 615264 &             146^2 & 127680 & 127680 & 231693 \\
%  34^2 & 3885 & 73920 & 73920 &                151^2 & 48488 & 128520 & 128520 \\
%  46^2 & 3360 & 3360 & 7293 &                  155^2 & 492936 & 546000 & 911064 \\
%  46^2 & 5280 & 9765 & 12768 &                 158^2 & 87261 & 184800 & 479136 \\
%  47^2 & 13167 & 349440 & 526320 &             162^2 & 155232 & 246240 & 279565 \\
%  47^2 & 122400 & 122400 & 401016 &            167^2 & 201552 & 425040 & 425040 \\
%  53^2 & 21216 & 28875 & 153216 &              168^2 & 418000 & 677376 & 981801 \\
%  53^2 & 488592 & 856520 & 875160 &            178^2 & 57120 & 200640 & 645645 \\
%  58^2 & 44160 & 94605 & 110880 &              191^2 & 354144 & 354144 & 550275 \\
%  59^2 & 25080 & 65688 & 93240 &               197^2 & 31416 & 73416 & 85800 \\
%  59^2 & 141680 & 575640 & 951048 &            201^2 & 53235 & 112480 & 399168 \\
%  62^2 & 6765 & 43680 & 43680 &                201^2 & 170280 & 294840 & 476560 \\
%  62^2 & 7392 & 16320 & 45885 &                213^2 & 19656 & 19656 & 68200 \\
% 62^2 & 8925 & 23712 & 36960 &                226^2 & 322245 & 350880 & 523488 \\
%  64^2 & 62985 & 270480 & 346368 &             233^2 & 39960 & 234080 & 923832 \\
%  67^2 & 47040 & 86112 & 126555 &              235^2 & 83904 & 96096 & 112875 \\
%  67^2 & 52632 & 64680 & 503880 &              238^2 & 47040 & 74400 & 460317 \\
%  67^2 & 89760 & 120120 & 146832 &             282^2 & 78880 & 125685 & 218592 \\
%  69^2 & 51408 & 196840 & 344520 &             290^2 & 686784 & 958341 & 976800 \\
%  70^2 & 36309 & 79200 & 114816 &              305^2 & 38019 & 317856 & 571200 \\
%  73^2 & 29640 & 224112 & 304920 &             331^2  & 429195 & 702240 & 896448 \\
%  73^2 & 31920 & 31920 & 123552 &              337^2  & 84456 & 231000 & 456456 \\
%  79^2 & 286440 & 286440 & 901968 &            358^2  & 192192 & 515040 & 780045 \\
%  83^2 & 451440 & 515840 & 832167 &            503^2  & 127680 & 127680 & 262515 \\
%  93^2 & 200200 & 407376 & 488376 &            573^2  & 140896 & 508896 & 675675 \\
%  100^2  & 309225 & 354816 & 836400 &           595^2  & 116571 & 501600 & 501600 \\
%  103^2  & 43680 & 43680 & 100947 &             605^2  & 253344 & 789600 & 993531 \\
 % 103^2  & 50895 & 118272 & 734160 &            879^2  & 99715 & 213408 & 993600 \\
%  103^2  & 105672 & 207480 & 333960 &           943^2  & 169592 & 582120 & 582120 \\
%  103^2  & 118272 & 312015 & 734160 &          &&&\\
% \hline
%\end{array}
%\end{equation*}
%}
\begin{table}[htbp]
\centering
\begin{tabular}{|c|c|c|c||c|c|c|c|}
\hline
$a_0$ & $a_1$ & $a_2$ & $a_3$ & $a_0$ & $a_1$ & $a_2$ & $a_3$ \\
\hline
 1 & 528 & 840 & 840 & $73^2$ & 31920 & 31920 & 123552 \\
 $2^2$ & 3360 & 3360 & 9405 & $93^2$ & 200200 & 407376 & 488376 \\
 $2^2$ & 4485 & 7392 & 20160 & $103^2$ & 43680 & 43680 & 100947 \\
 $2^2$ & 17952 & 26565 & 131040 & $103^2$ & 105672 & 207480 & 333960 \\
 $2^2$ & 100485 & 209760 & 376992 & $107^2$ & 18480 & 28152 & 127680 \\
 $5^2$ & 32736 & 46200 & 54264 & $107^2$ & 57195 & 364320 & 446880 \\
 $5^2$ & 78936 & 218064 & 319200 & $107^2$ & 102120 & 247632 & 414960 \\
 $7^2$ & 5280 & 5280 & 9555 & $109^2$ & 19803 & 36960 & 163680 \\
 $7^2$ & 40755 & 108192 & 252960 & $120^2$ & 106704 & 426496 & 482625 \\
 $7^2$ & 127400 & 180576 & 180576 & $125^2$ & 22011 & 84864 & 117600 \\
 $10^2$ & 2400 & 4389 & 8736 & $127^2$ & 33600 & 46371 & 280896 \\
 $13^2$ & 10440 & 15960 & 62832 & $129^2$ & 17955 & 40480 & 157248 \\
 $17^2$ & 45936 & 134400 & 249711 & $129^2$ & 177840 & 366520 & 506088 \\
 $23^2$ & 41496 & 138600 & 138600 & $130^2$ & 94656 & 181125 & 421344 \\
 $31^2$ & 10920 & 10920 & 26928 & $139^2$ & 115368 & 182280 & 318240 \\
 $31^2$ & 38640 & 57120 & 99528 & $141^2$ & 37240 & 201960 & 201960 \\
 $34^2$ & 3885 & 73920 & 73920 & $146^2$ & 127680 & 127680 & 231693 \\
 $46^2$ & 3360 & 3360 & 7293 & $151^2$ & 48488 & 128520 & 128520 \\
 $46^2$ & 5280 & 9765 & 12768 & $158^2$ & 87261 & 184800 & 479136 \\
 $47^2$ & 122400 & 122400 & 401016 & $162^2$ & 155232 & 246240 & 279565 \\
 $53^2$ & 21216 & 28875 & 153216 & $167^2$ & 201552 & 425040 & 425040 \\
 $58^2$ & 44160 & 94605 & 110880 & $197^2$ & 31416 & 73416 & 85800 \\
 $59^2$ & 25080 & 65688 & 93240 & $201^2$ & 53235 & 112480 & 399168 \\
 $62^2$ & 6765 & 43680 & 43680 & $201^2$ & 170280 & 294840 & 476560 \\
 $62^2$ & 7392 & 16320 & 45885 & $213^2$ & 19656 & 19656 & 68200 \\
 $62^2$ & 8925 & 23712 & 36960 & $226^2$ & 322245 & 350880 & 523488 \\
 $64^2$ & 62985 & 270480 & 346368 & $235^2$ & 83904 & 96096 & 112875 \\
 $67^2$ & 47040 & 86112 & 126555 & $238^2$ & 47040 & 74400 & 460317 \\
 $67^2$ & 52632 & 64680 & 503880 & $282^2$ & 78880 & 125685 & 218592 \\
 $67^2$ & 89760 & 120120 & 146832 & $337^2$ & 84456 & 231000 & 456456 \\
 $69^2$ & 51408 & 196840 & 344520 & $503^2$ & 127680 & 127680 & 262515 \\
 $70^2$ & 36309 & 79200 & 114816 & $595^2$ & 116571 & 501600 & 501600 \\
 $73^2$ & 29640 & 224112 & 304920 &  &  &  &  \\
 \hline
\end{tabular}
\caption{Reduced Hilbert cubes $H(a_{0}; a_{1}, a_{2}, a_{3})\subset \cal{S}$ with $\op{max}\{a_{0}, a_{1}, a_{2}, a_{3}\}\leq 2^{19}.$ }  
\end{table}
\bigskip
Our computations strongly suggest that there should be infinitely many reduced Hilbert cubes of dimension three contained in the set of squares. Moreover, among the solutions there are many solutions satisfying one of the additional conditions $a_{1}=a_{2}$ or $a_{2}=a_{3}$. This also suggests that there are infinitely many elements of $\cal{H}$ with only three different entries. 

As we will see in the next section, these expectations are true.

\section{Proof of Theorem \ref{thm:H3} and
Theorem  \ref{thm:idenical base}. }

\subsection{Construction of parametric solution of 3-dimensional cubes}
Let $a_{0}, a_{1}, a_{2}, a_{3}\in\Z$ and suppose that
$$
H(a_{0}; a_{1}, a_{2}, a_{3})=a_{0}+\{0, a_{1}\}+\{0, a_{2}\}+\{0, a_{3}\}\subset \cal{S}.
$$
We focus on the system (\ref{Hsystem}) and solve it by solving the equations one by one.

The first three equations can be easily solved by taking
$$
a_{0}=p^2, \quad a_{1}=q^2-p^2,\quad a_{2}=r^2-p^2.
$$
The fourth equation is equivalent with the equation $q^2-p^2=s^2-r^2$, and thus
$$
p=\frac{1}{2} (t z-x y),\quad q=\frac{1}{2} (t z+x y),\quad r=\frac{1}{2} (y z-t x),\quad s=\frac{1}{2} (t x+y z).
$$
We solve the fifth equation with respect to $a_{3}$ and get
$$
a_{3}=P^2-p^2=\frac{1}{4} \left(4 P^2-t^2 z^2+2 t x y z-x^2 y^2\right).
$$
The sixth equation takes the form
$$
P^2 + t x y z = Q^2
$$
and solving it with respect to $t$ we get $t=\frac{Q^2-P^2}{x y z}$. Substituting the computed values of $p, q, r, s, t$,
we get the corresponding values of $a_0$, $a_1$, $a_2$, $a_3$ as follows:
\begin{equation}\label{a0a1a2a3}
\begin{aligned}
	a_0 & = (P^2-Q^2+x^2y^2)^2/(4 x^2 y^2), \\
	a_1 & = Q^2-P^2, \\
	a_2 & = (x^2-z^2)(P^2-Q^2+x y^2 z)(P^2-Q^2-x y^2 z)/(4 x^2 y^2 z^2), \\
	a_3 & = -(P+Q-x y)(P+Q+x y)(P-Q+x y)(P-Q-x y)/(4 x^2y^2).
\end{aligned}
\end{equation}
Replacing 
$(R,S)$ by $(R,S)/(2 x y z)$, the remaining two equations take the form
\begin{equation}\label{sysRS}
\begin{aligned}
	R^2 & = x^2(P^2-Q^2)^2 + x^2y^4z^4 - ((P^2-Q^2)^2-x^2 y^2(2P-x y)(2P+x y
)) z^2, \\
	S^2 & = x^2(P^2-Q^2)^2 + x^2y^4z^4 - ((P^2-Q^2)^2-x^2 y^2(2Q-x y)(2Q+x y
)) z^2.
\end{aligned}
\end{equation}
Set $(P,Q) = (\frac{1}{2}(u y-m), \frac{1}{2}(u y+m))$ to give
\begin{equation}\label{RSeq}
\begin{aligned}
R^2 & = ((x^2-z^2)u^2 + x^2z^2)m^2 - 2u x^2z^2 m y + x^2z^2(u^2-x^2+z^2)y^2, \\ 
	S^2 & = ((x^2-z^2)u^2 + x^2z^2)m^2 + 2u x^2z^2 m y + x^2z^2(u^2-x^2+z^2)
y^2.
\end{aligned}
\end{equation}
This intersection in projective $m,y,R,S$-space is elliptic because of the point
\[ (m,y,R,S) = (x z, \; u, \; x z(x z-u^2), \; x z(x z+u^2)), \]
and a cubic model is 
\begin{align*}
	E: \;Y^2 =  X(X^2 + & 2( (x^2-z^2)u^4 - (x^4-4x^2z^2+z^4)u^2 -x^2z^2(x^2
-z^2)) X \\
	& + (u^2-x^2)^2 (x^2-z^2)^2 (u^2+z^2)^2).
\end{align*}
The curve $E$ has at least two independent points of infinite order,
\begin{align*}
	Q_1 & = ((x^2-u^2)^2z^2, \; (u^2-x^2)^2(u^2x^2+z^4)z), \\
	Q_2 & = ((x^2-z^2)^2u^2, \; (x^2-z^2)^2(u^4+x^2z^2)u).
\end{align*}
We can now pull back points $i_1 Q_1+i_2 Q_2$ to give parametrizations for $(a_0
,a_1,a_2,a_3)$ in terms of $u,x,z$. For example, $Q_2$ pulls back to
\[ (a_0,a_1,a_2,a_3) = (u^2(x+z)^2, \; -4u^2x z, \; 0, \; (u^2-x^2)(u^2-z^2)); \]
and $Q_1$ pulls back to
\begin{align*}
	&	(a_0,a_1,a_2,a_3) = ( \\
	& 	(2 x^2 z^2 u^6 + x^2(x^3-x^2z-2x z^2-6z^3) u^5 -2 z(x^5-x^4z-3x^
3z^2-5x^2z^3+z^5) u^4 \\
	& + 2 z^2(x+z)(x^4-2x^3z-2x^2z^2-2x z^3+z^4) u^3 - 2z^3(x^5-5x^3z^2-3x^2
z^3 \\
& -x z^4+z^5) u^2 - z^6(6x^3+2x^2z+x z^2-z^3) u + 2x^3z^7)^2, \\
	&  -4u x z(u^2 x^2 - 2u x^2 z + 2x^2 z^2 - z^4)(u^2 x^2 - 2u^2 z^2 + 2u 
z^3 - z^4) \times \\
	& (-u^2 x^2 + 2u^3 z - 2u^2 z^2 + 2u z^3 - z^4)(u^3 x^2 - 2u^2 x^2 z + 2
u x^2 z^2 - 2x^2 z^3 + u z^4), \\
	&  4(u - z)z(x^2 - z^2)(-x^2 + u z)(u^2 x^2 - z^4) (u^2 - u z + z^2) \times \\
	& (u^3 x - 2u^2 x z + u x^2 z - u^2 z^2 + 2u x z^2 - x z^3)(u^3 x - 2u^2
 x z - u x^2 z + u^2 z^2 + 2u x z^2 - x z^3), \\
	&  (u^2 - x^2)(u - z)(u^2 x^2 - 2u x^2 z + 2u z^3 - z^4)(u^2 x^2 - 2u^2 
x z + 2u x z^2 - 2x z^3 + z^4) \times \\
	& (u^2 x^2 + 2u^2 x z - 2u x z^2 + 2x z^3 + z^4)(u^3 x^2 - u^2 x^2 z + 2
u x^2 z^2 - 2u^2 z^3 + u z^4 - z^5)).
\end{align*}
The curve $E$ has a torsion point 
\[ Q_0 = ( -(u^2-x^2)(x^2-z^2)(u^2+z^2), \; 2u x z(x^2-z^2)(u^2-x^2)(u^2+z^2)) \]
of order 4, but the pullbacks of $i_0 Q_0+i_1 Q_1+i_2 Q_2$ give symmetries of the pullbacks with $i_0=0$, and so it is not necessary to consider the point $Q_0$.\\ \\
\noindent
We note that the curve (\ref{RSeq}) becomes obviously singular for $u^2-x^2+z^2=
0$, when it reduces to the genus 0 curve 
\begin{equation}\label{RSeq2}
\begin{aligned}
    R^2 & = m \, ((u^4 + x^2z^2)m - 2u x^2z^2 y), \\
	S^2 & = m \, ((u^4 + x^2z^2)m + 2u x^2z^2 y).
\end{aligned}
\end{equation}
Setting $(u,x,z)=(c^2-d^2,c^2+d^2,2c d)$, then
\begin{align*}
R^2 & = m( (c^8+14c^4d^4+d^8)m - 8c^2d^2(c^2-d^2)(c^2+d^2)^2y),\\
S^2 & = m( (c^8+14c^4d^4+d^8)m + 8c^2d^2(c^2-d^2)(c^2+d^2)^2y).
\end{align*}
From
\begin{align*}
& ( (c^8+14c^4d^4+d^8)m - 8c^2d^2(c^2-d^2)(c^2+d^2)^2y) \; \times \\
& ( (c^8+14c^4d^4+d^8)m + 8c^2d^2(c^2-d^2)(c^2+d^2)^2y)=\square
\end{align*}
there follows the parametrization 
\[ (m,y) = ( \frac{-A^2-B^2}{c^8+14c^4d^4+d^8}, \; \frac{A B}{4c^2d^2(c^2-d^2)
(c^2+d^2)^2} ), \]
% { (-A^2-B^2)/(c^8+14c^4d^4+d^8), A B/(4c^2d^2(c^2-d^2)(c^2+d^2)^2) }
with 
\[ (c^8+14c^4d^4+d^8) R^2 = (A+B)^2 (A^2+B^2), \quad (c^8+14c^4d^4+d^8)S^2 = (A-
B)^2 (A^2+B^2). \]
The latter are parameterized by
\begin{align*}
	(A,B,R,S) = ( -(c^4-d^4) G^2 - 8c^2d^2 G H + (c^4-d^4) H^2, & \\
    4 c^2d^2 G^2 - 2(c^4-d^4) G H - 4c^2d^2 H^2, & \\
	-(G^2 + H^2)((c^4-4c^2d^2-d^4)G^2 + 2(c^4+4c^2d^2-d^4) G H  - (c^4-4c^2d
^2-d^4) H^2), & \\
	(G^2 + H^2)((c^4+4c^2d^2-d^4)G^2 - 2(c^4-4c^2d^2-d^4) G H - (c^4+4c^2d^2
-d^4) H^2) & ).
\end{align*}
Then
\begin{align*}
	& (P,Q) = (\frac{u y-m}{2}, \; \frac{u y+m}{2}) = 1/(4c^2d^2(c^2+d^2)^2)
 \; \times \\
	& ( (4c^2d^4(c^2+d^2)G^4 + (c^8-18c^4d^4+d^8) G^3 H + 8c^2d^2(2c^2-d^2)(
c^2+d^2) G^2 H^2 \\
	& - (c^8-18 c^4d^4+d^8) G H^3 + 4c^2d^4(c^2+d^2) H^4), \\
	& (-4c^4d^2(c^2+d^2)G^4 + (c^8-18c^4d^4+d^8) G^3 H + 8c^2d^2(c^2-2d^2)(c
^2+d^2) G^2 H^2 \\
	& -(c^8-18c^4d^4+d^8) G H^3 -4 c^4d^2(c^2+d^2) H^4)).
\end{align*}
This gives
\[ t=(Q^2-P^2)/(x y z) = -(c^2-d^2)(G^2+H^2)^2/(2c d(c^2+d^2)) \]
% -(c^2-d^2)*(G^2+H^2)^2/(2*c*d*(c^2+d^2))
from which $(p,q,r,s)=\left( \frac{t z-x y}{2}, \frac{t z+x y}{2}, \frac{y z-t x
}{2}, \frac{y z+t x}{2} \right)$
may be computed, leading to \\
\begin{equation}\label{a1a2a3a4solution}
\begin{aligned}
	& a_0 \; : \; a_1 \; : \; a_2 \; : \; a_3 = p^2 \; : \; q^2-p^2 \; : \; 
r^2-p^2 \; : \; P^2-p^2 = \\
	& \\
& (c^2+d^2)^2 (-4c^2d^4(c^2-d^2) G^4 + (c^8-18c^4d^4+d^8) G^3 H \\
& + 8c^2d^2(c^2-d^2)(2c^2+d^2) G^2 H^2 - (c^8-18c^4d^4+d^8) G H^3 - 4c^2d^4(c^2-d^2) H^4)^2: \\
	& \\
& 8c^2d^2(c^4-d^4)^2 (G^2+H^2)^2 (2c^2d^2 G^2 - (c^4-d^4) G H - 2c^2d^2 H^2)((c^
4-d^4) G^2 \\
& + 8c^2d^2 G H - (c^4-d^4)H^2): \\
	& \\
& (c^2-d^2)^2(d(c+d) G - c(c-d) H)((d(c-d) G + c(c+d) H) \times \\
	& ((c^2-2c d-d^2) G + (c^2+d^2) H) ((c^2+2c d-d^2) G + (c^2+d^2) H) \\
	& ((c^2+d^2) G - (c^2+2c d-d^2) H) (c(c+d) G - d(c-d) H) \times\\
	& (c(c-d) G + d(c+d) H)((c^2+d^2) G - (c^2-2 c d-d^2) H)): \\
	& \\
& 4c^2d^2G H(G^2-H^2)((c^4-d^4) G + 4c^2d^2 H)(4c^2d^2 G - (c^4-d^4) H) \times \\
	& ((c^4-4c^2d^2-d^4) G + (c^4+4c^2d^2-d^4) H) ((c^4+4c^2d^2-d^4) G - (c^
4-4c^2d^2-d^4) H).
\end{aligned}
\end{equation}
Note that $a_0$, $a_1$, $a_2$, $a_3$ treated as polynomials in $\Z[c,d,G,H]$
are homogeneous with respect to $c,d$ of degree $20$,  while they are homogeneous with respect to $G,H$ of degree $8$. We are in the position to present the lower bound for $H_{3}(N)$.


\begin{proof}[Proof of Theorem \ref{thm:H3}] Here is the explanation how to get $\gg N^{1/8}$ reduced Hilbert cubes in $[0, N]$. We take $(c,d)=(3,1)$ in (\ref{a1a2a3a4solution}) and replace $(G,H)$ by $(t,1)$ in the parametrization (\ref{a1a2a3a4solution}). Next, applying the map $\phi_{1}$ and then the map $\phi_{\sigma}$ which switches $a_{1}$ with $a_{3}$, we get the Hilbert cube $H(a_{0}; a_{1}, a_{2}, a_{3})$, with
\begin{align*}
a_{0}&=25 \left(18 t^4-319 t^3-684 t^2+319 t+18\right)^2 \\
a_{1}&=9 (t-1) t (t+1) (9 t-20) (11 t+29) (20 t+9) (29 t-11)\\
a_{2}&= 16 (t+5)(t+6)(2t-3)(3t+2)(5t-7)(5t-1) (6 t-1) (7 t+5) \\
a_{3}&=7200(t^2+1)^2(9 t^2-40 t-9)(10 t^2+9 t-10).  \\
\end{align*}
For $t\geq 7$ all entries are positive and satisfy $a_{1}<a_{2}<a_{3}$. Next, in order to have a reduced cube we note that the polynomials $a_{i}(t), i=0, 1, 2, 3$ are co-prime as polynomials in $\Z[t]$. This means that there exists an integer $M$ such that a possible common factor is bounded independently of $t$. To be more precise, it is enough to take 
$$
M=2^{45}\cdot 3^{18}\cdot 5^{2}\cdot 13^{16}\cdot 37^{16}.
$$
The shape of this number follows from the extended Euclidean algorithm applied to the polynomials $a_{0},a_{1}$. 

Because the number of possible common divisors is finite, the number of reduced Hilbert cubes grows like $N^{1/8}$ with $N$ going to infinity.
\end{proof}




%\section{Small solutions and related observations}

%To look for small Hilbert cubes of dimension 3 sitting in the set of squares one can %look for the solutions of system (\ref{sysRS}). With each solution of (\ref{sysRS}) we %have the corresponding Hilbert cube $H(a_{0}; a_{1}, a_{2}, a_{3})$. However, many %different values of parameters $x, y, z, P, Q$ may lead to the same Hilbert cube $H$. % Taking into account this observation, and noting that we can assume that $x, y, z, P, %Q$ are positive (see expressions (\ref{a0a1a2a3})), we performed small numerical %search for integer solutions of the system (\ref{sysRS}) with $\max\{x, y, z, P, %Q\}\leq 100$ and such that the corresponding values of $a_{i}, i=0, 1, 2, 3$ are %non-zero. In the considered range we found 50 non-trivial solutions. We implemented %the procedure search in a simple {\tt Pari} script \cite{PARI2}. The computations were %performed on a typical laptop with i7 type processor of 9th generation and 32 GB of %memory and took 2 h 10 min. We present results of our computations in the table below. 
%\begin{equation*}
%\begin{array}{|ccccc||llll|}
%\hline
% x & y & z & P & Q  & a_0 & a_1 & a_2 & a_3 \\
% \hline
% 1 & 22 & 3 & 29 & 37 & 1 & 528 & 840 & 840 \\
% 2 & 15 & 7 & 29 & 41 & 4 & 3360 & 3360 & 9405 \\
% 1 & 44 & 2 & 71 & 83 & 4 & 4485 & 7392 & 20160 \\
% 3 & 11 & 8 & 49 & 61 & 49 & 5280 & 5280 & 9555 \\
% 1 & 40 & 5 & 67 & 83 & 100 & 2400 & 4389 & 8736 \\
% 5 & 7 & 3 & 50 & 62 & 289 & 45936 & 134400 & 249711 \\
% 5 & 7 & 33 & 54 & 65 & 324 & 294525 & 655776 & 2855776 \\
% 7 & 10 & 18 & 41 & 85 & 529 & 41496 & 138600 & 138600 \\
% 5 & 8 & 21 & 71 & 83 & 961 & 184800 & 503139 & 615264 \\
% 2 & 7 & 15 & 37 & 47 & 2116 & 3360 & 3360 & 7293 \\
% 1 & 20 & 6 & 61 & 71 & 2116 & 5280 & 9765 & 12768 \\
% 13 & 2 & 20 & 23 & 35 & 2500 & 2155461 & 8937600 & 11762400 \\
% 5 & 5 & 3 & 79 & 86 & 2809 & 21216 & 28875 & 153216 \\
% 2 & 19 & 5 & 59 & 71 & 3364 & 2225685 & 2252640 & 5023200 \\
% 1 & 40 & 6 & 53 & 83 & 3844 & 7392 & 16320 & 45885 \\
% 3 & 28 & 1 & 37 & 83 & 4096 & 62985 & 270480 & 346368 \\
% 3 & 14 & 8 & 43 & 53 & 4489 & 47040 & 86112 & 126555 \\
% 2 & 22 & 5 & 17 & 53 & 5329 & 29640 & 224112 & 304920 \\
% 1 & 32 & 7 & 31 & 53 & 10609 & 50895 & 118272 & 734160 \\
% 2 & 22 & 5 & 31 & 61 & 10609 & 105672 & 207480 & 333960 \\
% 1 & 32 & 7 & 71 & 83 & 10609 & 118272 & 312015 & 734160 \\
% 8 & 7 & 19 & 23 & 53 & 11449 & 92235 & 446880 & 646272 \\
% 8 & 7 & 47 & 35 & 59 & 12100 & 228000 & 442176 & 5072925 \\
% 5 & 7 & 16 & 25 & 31 & 16129 & 33600 & 46371 & 280896 \\
% 7 & 13 & 9 & 31 & 73 & 16641 & 17955 & 40480 & 157248 \\
% 11 & 6 & 6 & 49 & 93 & 16641 & 177840 & 366520 & 506088 \\
% 1 & 32 & 7 & 67 & 81 & 17161 & 132608 & 270135 & 720720 \\
% 10 & 7 & 27 & 75 & 93 & 17956 & 302400 & 544544 & 730269 \\
% 10 & 3 & 7 & 71 & 83 & 24964 & 87261 & 184800 & 479136 \\
% 1 & 84 & 2 & 79 & 97 & 26244 & 155232 & 246240 & 279565 \\
% 5 & 5 & 22 & 42 & 53 & 28224 & 418000 & 677376 & 981801 \\
% 5 & 22 & 1 & 53 & 79 & 31416 & 38809 & 73416 & 85800 \\
% 
 %\hline
%\end{array}
%\end{equation*}
%\begin{equation*}
%\begin{array}{|ccccc||llll|}
%\hline
% x & y & z & P & Q  & a_0 & a_1 & a_2 & a_3 \\
% \hline
% 16 & 3 & 21 & 17 & 39 & 40401 & 53235 & 112480 & 399168 \\
% 7 & 6 & 30 & 51 & 79 & 40401 & 170280 & 294840 & 476560 \\
% 2 & 78 & 3 & 85 & 97 & 19656 & 19656 & 45369 & 68200 \\
% 45 & 2 & 88 & 41 & 85 & 80656 & 591744 & 2007369 & 2217600 \\
% 13 & 5 & 24 & 49 & 91 & 109561 & 742368 & 1513515 & 3974880 \\
% 13 & 4 & 7 & 29 & 63 & 137641 & 6826680 & 14370840 & 25902968 \\
% 13 & 5 & 3 & 23 & 53 & 151321 & 206283 & 1541280 & 3020640 \\
% 95 & 1 & 33 & 58 & 96 & 251001 & 2776599 & 4436224 & 5266800 \\
% 13 & 9 & 32 & 42 & 94 & 259081 & 312455 & 2291328 & 5362560 \\
% 35 & 3 & 39 & 61 & 82 & 140896 & 328329 & 508896 & 675675 \\
% 2 & 19 & 5 & 61 & 81 & 487204 & 1042965 & 4100960 & 4885920 \\
% 5 & 5 & 69 & 87 & 97 & 531441 & 1656000 & 6280659 & 25427584 \\
% 4 & 11 & 23 & 51 & 87 & 574564 & 684320 & 2404512 & 6014925 \\
% 1 & 9 & 10 & 64 & 69 & 846945 & 861840 & 1364224 & 3944192 \\
% 3 & 7 & 20 & 37 & 52 & 1420864 & 2872320 & 4186560 & 11923545 \\
% 8 & 13 & 33 & 35 & 87 & 2013561 & 3322539 & 27634464 & 184828000 \\
% 35 & 3 & 69 & 71 & 90 & 2913849 & 6194475 & 7294176 & 15791776 \\
% 1 & 9 & 33 & 64 & 67 & 2944656 & 3851793 & 37200240 & 211267840 \\
% \hline
%\end{array}
%\end{equation*}
%\begin{center} Table 1. Some solutions of the system (\ref{sysRS}) and the %corresponding Hilbert cubes\end{center}


\subsection{Hilbert cubes with two identical base elements}
\bigskip

From Table 1 we see that there are also Hilbert cubes $H(a_{0}; a_{1}, a_{2}, a_{3})$ satisfying the condition $a_{1}=a_{2}$. We prove that there are infinitely many such examples.

\begin{proof}[Proof of Theorem \ref{thm:idenical base}]
The system guaranteeing that $H(a_{0}; a_{1}, a_{2}, a_{3})\in\cal{S}$ with $a_{1}=a_{2}$ takes the form
\begin{center}
\begin{tabular}{rrr}
        $a_0=p^2$, & $a_0+a_1=q^2$, & $a_0+2a_1=r^2$, \\
        $a_0+a_3=s^2$, & $a_0+a_1+a_3=t^2$, & $a_0+2a_1+a_3=u^2$,
\end{tabular}
\end{center}
so that $(a_0,a_1,a_3) = (p^2, q^2-p^2, s^2-p^2)$ with
\[ p^2-2q^2+r^2 = 0, \; p^2-q^2-s^2+t^2 = 0, \; 2p^2-2q^2-s^2+u^2 = 0. \]
Set
\[ (p,q,s,t) = (a d-b c, \; a c-b d, \; a d+b c, \; a c+b d) \]
to give
\begin{align*}
(2c^2-d^2) a^2 - 2c d a b - (c^2-2d^2) b^2 & = r^2, \\
(2c^2-d^2) a^2 + 2c d a b - (c^2-2d^2) b^2 & = u^2
\end{align*}
a curve of genus 1 over the function field $\Q(c,d)$,
and in fact elliptic since there is a point at $(a,b,r,u)=(1,1,c-d,c+d)$.
A cubic model is
\[ Y^2 = X ( X^2 - 4(c^2-c d-d^2)(c^2+c d-d^2) X + 4(c^2-d^2)^4), \]
with points
\[ P_0(X,Y)=((c^2-d^2)^2, \; (c^2-d^2)^2(c^2+d^2)), \quad P_1(X,Y)=(2(c^2-d^2)^2, \; 4c d(c^2-d^2)^2) \]
of infinite order and order 4, respectively.
The pullback of multiples of $P_0$ give solutions $(a_0,a_1,a_2,a_3)$ for which $a_1=a_2$. For example, the point $2P_0$ pulls back to
\begin{align*}
        (a,b,r,u) = ( & (5c^2-d^2)(c^2+7d^2), \; -(c^2-5d^2)(7c^2+d^2), \\
        & (c-d)(c^4+36c^3d+38c^2d^2+36c d^3+d^4), \\
        & (c+d)(c^4-36c^3d+38c^2d^2-36c d^3+d^4))
\end{align*}
with
\begin{equation}\label{a1a2solution}
\begin{aligned}
        a_0 & = (c-d)^2(7c^4+12c^3d-22c^2d^2+12c d^3+7d^4)^2, \\
        a_1 & = -24(c-d)^2(c+d)^2(c^2+d^2)(c^2-6c d+d^2)(c^2+6c d+d^2), \\
        a_3 & = -4c d(c^2-5d^2)(5c^2-d^2)(7c^2+d^2)(c^2+7d^2)
\end{aligned}
\end{equation}

\end{proof}
\begin{rem}
{\rm We can also find reduced Hilbert cubes $H(a'_{0};a'_{1},a'_{2},a'_{3})$ where $a'_2=a'_3$. 
Indeed, it is enough at (\ref{a1a2solution}) to find conditions for $c,d$ such that $a_{i}(c,d)>0$ and $a_{2}(c,d)>a_{3}(c,d)$. Then $H(a_{0};a_{3},a_{1},a_{2})$ is the Hilbert cube we are looking for. Because the polynomials $a_{i}(c,d)$ are homogeneous we can assume that $d=1$ and look for solutions of the corresponding inequalities for rational $c$.
One can easily check that if 
$$
c \in\left(\frac{1}{4} \left(3+\sqrt{17}\right), \sqrt{5} \right)=:I,
$$
then $a_{i}(c,1)>0$ and $a_{2}(c,1)>a_{3}(c,1)$. Because there are infinitely many rational numbers in $I$, we get infinitely many Hilbert cubes satisfying the required condition. 
}    
\end{rem}
\bigskip

\section{Proof of corollary {\ref{cor-fractions}} }

We apply our parametric families of Hilbert cubes contained in the set of squares 
to prove the corollary. We restate it for convenience:

Let $\cal{H}_{+}=\cal{H}\cap \N_{+}^{4}$. For each $i, j\in\{0, 1, 2, 3\}$ with $i<j$, the set
$$
H_{i,j}:=\left\{\frac{a_{i}}{a_{j}}:\;(a_{0}, a_{1}, a_{2}, a_{3})\in\cal{H}_{+}\right\}\subset \Q_{+}
$$
is dense in the Euclidean topology in the set $\R_{+}$.

\begin{proof}
The proof in each case is similar, and we prove the result for the cases $i < j$.
We first show $\overline{H_{0,1}}=\R_{+}$. Define the function
\[ f_{0,1}(x) = \frac{a_0(x,1)}{a_1(x,1)} = \frac{ -(7+12x-22x^2+12x^3+7x^4)^2} {24(1 + x)^2(1 + x^2)(1-6x+x^2)(1+6x+x^2)} \]
where $a_0$, $a_1$ are the polynomials from (\ref{a1a2solution}). Note that $f_{0,1}(c/d) = a_0(c,d)/a_1(c,d)$. \\
The functions $a_i(x,1)$ for $i=1,2,3$, are all positive on the intervals
\[ x \in (-3-2\sqrt{2}, \; -\sqrt{5}) \; \cup \; \left(-\frac{1}{\sqrt{5}}, \; -\frac{3+\sqrt{17}}{4}\right) \; \cup \; \left(\frac{1}{\sqrt{5}}, \; \sqrt{5}\right), \]
so that $a_0(x,1)$ being square implies $f_{0,1}(x)$ is non-negative on these ranges.
The function $f_{0,1}(x)$ has a local minimum value of $0$ on the interval $(-3-2\sqrt{2}, \; -\sqrt{5})$, occurring at $x_{min} = \frac{(-3-2\sqrt[4]{2})}{2\sqrt{2}-1} \sim -2.94155$; define $I_{0,1} = (-3-2\sqrt{2}, \; x_{min})$.
We have
\[ \lim_{x \rightarrow (-3-2\sqrt{2})^{+}} f_{0,1}(x) = +\infty, \]
and $f_{0,1}(x)$ is continuous and decreasing on $(-3-2\sqrt{2}, \; x_{min})$.
It follows that
\[ \{ f_{0,1}(x) : x \in \Q_{+} \cap I_{0,1} \} \subset H_{0,1}. \]
The density of $\Q_{+} \cap I_{0,1}$ in $I_{0,1}$ implies that the Euclidean closure of the set $H_{0,1}$ is dense in $\R_{+}$.
Because we are in the situation $a_1=a_2$ we also get that $H_{0,2}$ is dense in $\R_{+}$. \\ \\
For the density of $H_{0,3}$, we follow the same approach and define
\[ f_{0,3}(x) = \frac{a_0(x,1)}{a_3(x,1)} = \frac{-(x-1)^2 (7+12x-22x^2+12x^3+7x^4)^2}{4x(x^2-5)(5x^2-1)(x^2+7)(7x^2+1)}. \]

Let $I_{0,3}$ be the interval $(1, \sqrt{5})$, on which $a_i(x,1)$, $i=0,...,3$, are all positive. Then
\[ \{ f_{0,3}(x): x \in \Q_{+} \cap I_{0,3} \} \subset H_{0,3}. \]
Because $f_{0,3}$ is continuous and increasing on $I_{0,3}$, and $f_{0,3}(1)=0$, and
\[ \lim_{x \rightarrow \sqrt{5}^{-}} f_{0,3}(x) = +\infty, \]
it follows that $H_{0,3}$ is dense in $\R_{+}$. \\ \\
To prove density of $H_{1,3}$, consider the function
\[ f_{1,3}(x) = \frac{a_1(x,1)}{a_3(x,1)} = \frac{ (6(1-x)^2(1+x)^2(1+x^2)(1-6x+x^2)(1+6x+x^2))}{ x(5-x^2)(1-5x^2)(7+x^2)(1+7x^2)}. \]
Let $I_{1,3} = (1, \sqrt{5})$, on which $a_i(x,1)$, $i=0,...,3$, are all positive. Then
\[ \{ f_{1,3}(x): x \in \Q_{+} \cap I_{1,3} \} \subset H_{1,3}. \]
Because $f_{1,3}$ is continuous and increasing on $I_{1,3}$, and $f_{1,3}(1)=0$, and
\[ \lim_{x \rightarrow \sqrt{5}^{-}} f_{1,3}(x) = +\infty, \]
it follows that $H_{1,3}$ is dense in $\R_{+}$. By applying the same function, we deduce also that $H_{2,3}$ is dense in $\R_{+}$. \\ \\
Finally, to get the density of $H_{1,2}$ we cannot of course use the polynomials (\ref{a1a2solution}), where $a_1=a_2$. Instead, we need to go back to the parametric solution given by (\ref{a1a2a3a4solution}) and consider for example the function
\begin{align*}
  f_{1,2}(x) & = \frac{a_1(x,1,2,1)}{a_2(x,1,2,1)} \\
  & = \frac{400x^2(1+x^2)^2(-1-3x^2+x^4)(-3+16x^2+3x^4)}{(9+2x^2+x^4)(4-13x^2+x^4)(1-22x^2+9x^4)(1+3x^2+4x^4)}.
\end{align*}
Let $I_{1,2}=\left( \frac{2+\sqrt{7}}{3}, \sqrt{ \frac{3+\sqrt{13}}{2}} \right)$, on which interval $a_i(x,1,2,1)$, $i=0,1,2,3$, are all positive. Then
\[ \{ f_{1,2}(x): x \in \Q_{+} \cap I_{1,2} \} \subset H_{1,2}. \]
Since $f_{1,2}$ is continuous and decreasing on $I_{1,2}$, and
$f_{1,2}\left( \sqrt{ \frac{3+\sqrt{13}}{2}} \right)=0$, and
\[ \lim_{x \rightarrow \left( \frac{2+\sqrt{7}}{3} \right)^{+}} f(1,2)(u) = +\infty, \]
it follows that $H_{1,2}$ is dense in $\R_{+}$.

\end{proof}

\section{Cubes with the same values of $a_0$, $a_1$, $a_2$.}
We have observed numerically many instances of Hilbert cube pairs $(a_0,a_1,a_2,a_3)$, $(A_0,A_1,A_2,A_3)$, in which $a_i=A_i$, $i=0,1,2$. More precisely, in the range $a_{0}<a_{1}<10^8$ there are exactly 6 examples of this kind presented in the table below.
\begin{table}[htbp]
\centering
\begin{tabular}{|l|l|l|l|l|}
\hline
$a_{0}=A_{0}$ & $a_{1}=A_{1}$ & $a_{2}=A_{2}$ & $a_{3}$ & $A_{3}$ \\
\hline
332929  & 6726720  & 6726720  & 8322435   & 22381827 \\
438244  & 1004157  & 1939520  & 3013920   & 8791200  \\
643204  & 1367520  & 1367520  & 6804237   & 35947197 \\
4674244 & 1367520  & 1367520  & 2773197   & 31916157 \\
4713241 & 71831760 & 71831760 & 130613448 & 665527080 \\
38775529 & 71831760 & 71831760 & 96551160  & 631464792 \\
\hline
\end{tabular}
\caption{ Hilbert cubes $H(a_0; a_1, a_2, a_3), H(A_0; A_1, A_2, A_3)\subset \cal{S}$ with $a_{i}=A_{i}$ for $i=0, 1, 2$.}
\end{table}
 Based on this (modest) set of examples one can speculate that there should be infinitely many such examples. As our next result shows our expectation is true.
\begin{thm}
    There are infinitely many pairs $H(a_{0}; a_{1}, a_{2}, a_{3}), H(A_0; A_1, A_2, A_3)\subset \cal{S}$ of Hilbert subes such that $a_{i}=A_{i}$ for $i=0, 1, 2$.
\end{thm}
\begin{proof} 
The defining system is
\begin{center}
\begin{tabular}{rrrr}
	$a_0=p^2$ & $a_0+a_1=q^2$ & $a_0+a_2=r^2$ & $a_0+a_1+a_2+a_3=s^2$ \\
	$a_0+a_3=P_1^2$ & $a_0+a_1+a_3=Q_1^2$ & $a_0+a_2+a_3=R_1^2$ & $a_0+a_1+a_2+a_3=S_1^2$ \\
	$a_0+A_3=P_2^2$ & $a_0+a_1+A_3=Q_2^2$ & $a_0+a_2+A_3=R_2^2$ & $a_0+a_1+a_2+A_3=S_2^2$.
\end{tabular}
\end{center}
On eliminating $a_0,a_1,a_2,a_3,A_3$, this is equivalent to the system
\[ p^2+s^2=q^2+r^2, \qquad P_1^2+S_1^2=Q_1^2+R_1^2, \qquad P_2^2+S_2^2=Q_2^2+R_2^2, \]
\[ p^2-q^2 = P_1^2-Q_1^2 = P_2^2-Q_2^2, \qquad p^2-r^2 = P_1^2-R_1^2 = P_2^2-R_2^2. \]
To satisfy the first row, set
\begin{align*}
	(p,q,r,s) = & (\alpha_0 \delta_0-\beta_0 \gamma_0, \;\; \alpha_0 \delta_0+\beta_0 \gamma_0, \;\; \alpha_0 \gamma_0-\beta_0 \delta_0, \;\; \alpha_0 \gamma_0+\beta_0 \delta_0), \\
	(P_1,Q_1,R_1,S_1) = & (\alpha_1 \delta_1-\beta_1 \gamma_1, \;\; \alpha_1
 \delta_1+\beta_1 \gamma_1, \;\; \alpha_1 \gamma_1-\beta_1 \delta_1, \;\; \alpha_1 \gamma_1+\beta_1 \delta_1), \\
	(P_2,Q_2,R_2,S_2) = & (\alpha_2 \delta_2-\beta_2 \gamma_2, \;\; \alpha_2
 \delta_2+\beta_2 \gamma_2, \;\; \alpha_2 \gamma_2-\beta_2 \delta_2, \;\; \alpha_2 \gamma_2+\beta_2 \delta_2).
\end{align*}
The second row then delivers
\[ \alpha_0 \beta_0 \gamma_0 \delta_0 = \alpha_1 \beta_1 \gamma_1 \delta_1 = \alpha_2 \beta_2 \gamma_2 \delta_2, \]
\[ (\alpha_0^2-\beta_0^2)(\gamma_0^2-\delta_0^2) = (\alpha_1^2-\beta_1^2)(\gamma
_1^2-\delta_1^2) = (\alpha_2^2-\beta_2^2)(\gamma_2^2-\delta_2^2). \] 
If we set
\[ (\alpha_i,\beta_i)=\left( \frac{\gamma_i+\delta_i}{2},\;\frac{\gamma_i-\delta_i}{2} \right), \quad i=0,1,2, \]
then we obtain the simple system
\[ \gamma_0 \delta_0(\gamma_0^2-\delta_0^2) = \gamma_1 \delta_1(\gamma_1^2-\delta_1^2) = \gamma_2 \delta_2(\gamma_2^2-\delta_2^2), \]
requiring three Pythagorean triangles with equal area. Such are provided by
\begin{align*}
	(\gamma_0,\delta_0) = & (u^2+u v+v^2, u^2-v^2), \\
	(\gamma_1,\delta_1) = & (u^2+u v+v^2, 2u v+v^2), \\
	(\gamma_2,\delta_2) = & (u^2+2 u v, u^2+u v+v^2),
\end{align*}
with 
\begin{align*}
	(\alpha_0,\beta_0) = & (u(2u+v)/2, v(u+2v)/2), \\
	(\alpha_1,\beta_1) = & ((u+v)(u+2v)/2, u(u-v)/2) \\
	(\alpha_2,\beta_2) = & ((u+v)(2u+v)/2, (u-v)v/2).
\end{align*}
Then
\begin{align*}
	(p,q,r,s) = & ((2u^4 - 5u^2v^2 - 4uv^3 - 2v^4)/2, \\
	& (2u^4 + 2u^3v + u^2v^2 + 2uv^3 + 2v^4)/2, \\
	& (2u^4 + 2u^3v + u^2v^2 + 2uv^3 + 2v^4)/2, \\
	& (2u^4 + 4u^3v + 5u^2v^2 - 2v^4)/2);
\end{align*}
% {(2*u^4 - 5*u^2*v^2 - 4*u*v^3 - 2*v^4)/2, (2*u^4 + 2*u^3*v + u^2*v^2 + 2*u*v^3 + 2*v^4)/2, (2*u^4 + 2*u^3*v + u^2*v^2 + 2*u*v^3 + 2*v^4)/2, (2*u^4 + 4*u^3*v + 5*u^2*v^2 - 2*v^4)/2}
\begin{align*}
	(P_1,Q_1,R_1,S_1) = & ((-u^4 + 2u^3v + 7u^2v^2 + 8uv^3 + 2v^4)/2, \\
	& (u^4 + 2u^3v + 7u^2v^2 + 6uv^3 + 2v^4)/2,  \\
	& (u^4 + 2u^3v + 7u^2v^2 + 6uv^3 + 2v^4)/2, \\
	& (u^4 + 6u^3v + 5u^2v^2 + 4uv^3 + 2v^4)/2)
\end{align*}
% {(-u^4 + 2*u^3*v + 7*u^2*v^2 + 8*u*v^3 + 2*v^4)/2, (u^4 + 2*u^3*v + 7*u^2*v^2 + 6*u*v^3 + 2*v^4)/2, (u^4 + 2*u^3*v + 7*u^2*v^2 + 6*u*v^3 + 2*v^4)/2, (u^4 + 6*u^3*v + 5*u^2*v^2 + 4*u*v^3 + 2*v^4)/2}
\begin{align*}
	(P_2,Q_2,R_2,S_2) = & ((2u^4 + 4u^3v + 5u^2v^2 + 6uv^3 + v^4)/2, \\
	& (2u^4 + 6u^3v + 7u^2v^2 + 2uv^3 + v^4)/2, \\
	& (2u^4 + 6u^3v + 7u^2v^2 + 2uv^3 + v^4)/2, \\
	& (2u^4 + 8u^3v + 7u^2v^2 + 2uv^3 - v^4)/2).
\end{align*}
% {(2*u^4 + 4*u^3*v + 5*u^2*v^2 + 6*u*v^3 + v^4)/2, (2*u^4 + 6*u^3*v + 7*u^2*v^2 + 2*u*v^3 + v^4)/2, (2*u^4 + 6*u^3*v + 7*u^2*v^2 + 2*u*v^3 + v^4)/2, (2*u^4 + 8*u^3*v + 7*u^2*v^2 + 2*u*v^3 - v^4)/2}
Finally, scaling by a factor 4, 
\begin{align*}
	(a_0,a_1,a_2,a_3) = &  ( (2u^4 - 5u^2v^2 - 4uv^3 - 2v^4)^2, \\
	& 4u(u - v)v(u + v)(2u + v)(u + 2v)(u^2 + uv + v^2), \\
	& 4u(u - v)v(u + v)(2u + v)(u + 2v)(u^2 + uv + v^2), \\
	& -u(u + 2v)(u^2 - 2uv - 2v^2)(u^2 + 2v^2)(3u^2 + 4uv + 2v^2) );
\end{align*}
% {(2*u^4 - 5*u^2*v^2 - 4*u*v^3 - 2*v^4)^2/4, u*(u - v)*v*(u + v)*(2*u + v)* (u+ 2*v)*(u^2 + u*v + v^2), u*(u - v)*v*(u + v)*(2*u + v)*(u + 2*v)* (u^2 + u*v + v^2), -1/4*(u*(u + 2*v)*(u^2 - 2*u*v - 2*v^2)*(u^2 + 2*v^2)* (3*u^2 + 4*u*v + 2*v^2))}
\begin{align*}
	(A_0,A_1,A_2,A_3) = & ((2u^4 - 5u^2v^2 - 4uv^3 - 2v^4)^2, \\
	& 4u(u - v)v(u + v)(2u + v)(u + 2v)(u^2 + uv + v^2), \\
	& 4u(u - v)v(u + v)(2u + v)(u + 2v)(u^2 + uv + v^2), \\
	& v(2u + v)(2u^2 + 2uv - v^2)(2u^2 + v^2)(2u^2 + 4uv + 3v^2)),
\end{align*}
% {(2*u^4 - 5*u^2*v^2 - 4*u*v^3 - 2*v^4)^2/4, u*(u - v)*v*(u + v)*(2*u + v)* (u+ 2*v)*(u^2 + u*v + v^2), u*(u - v)*v*(u + v)*(2*u + v)*(u + 2*v)* (u^2 + u*v +v^2), (v*(2*u + v)*(2*u^2 + 2*u*v - v^2)*(2*u^2 + v^2)* (2*u^2 + 4*u*v + 3*v^2))/4}
where $a_0=A_0$, $a_1=A_1$, $a_2=A_2$.
\end{proof}

\section{Conjectures, remarks, computational observations}\label{sec4}

In this section we formulate some conjectures which may stimulate further research. We already answered on the second and fourth part of Question \ref{genques}. We used our computational approach to study the first part of Question \ref{genques}. More precisely, we are interested whether for given $n\in\N$ there is a reduced Hilbert cube $H(a_{0}; a_{1}, a_{2}, a_{3})$ such that $a_{0}=n^2$. Note that the assumption that $H(a_{0}; a_{1}, a_{2}, a_{3})$ need to be reduced is important and makes the question nontrivial. Indeed, without this assumption we get the Hilbert cube $H(n^2, 528n^2, 840n^2, 840n^2)\subset \cal{S}$. 

Let us note that if $a_{0}=n^2$ is fixed then, necessarily $a_{1}=t^2-n^2$ and using small modification of Algorithm 1 we hunted for reduced Hilbert cubes of dimension three with $a_{0}=n^2$ and $a_{1}\leq (1.4\cdot 10^8)^2$ for $n\in\{1, \ldots, 100\}$. These computations took several days on the personal computerof the third author. In this range we found the appropriate solution for all but four numbers given by $n=42, 44, 56,  80$. For some values of $n$ the examples are quite large. We collect our findings in the table below. 
\begin{landscape}
\begin{table}[htbp]
\centering
\begin{tabular}{|l|l|l|l||l|l|l|l|}
\hline
$a_{0}$ & $a_{1}$ & $a_{2}$ & $a_{3}$ &
$a_{0}$ & $a_{1}$ & $a_{2}$ & $a_{3}$ \\
\hline
$1^2$  & 528 & 840 & 840 & $26^2$ & 151826342525 & 3489610801824 & 66625590001824 \\
$2^2$  & 3360 & 3360 & 9405 & $27^2$ & 8334040 & 15704640 & 33280632 \\
$3^2$  & 13704795 & 35772352 & 57562560 & $28^2$ & 44782080 & 97041417 & 730728240 \\
$4^2$  & 565488 & 3094065 & 4309760 & $29^2$ & 112728 & 159960 & 1584240 \\
$5^2$  & 32736 & 46200 & 54264 & $30^2$ & 21903703101 & 151484423200 & 356015088000 \\
$6^2$  & 7215822880 & 10763232480 & 26819630253 & $31^2$ & 10920 & 10920 & 26928 \\
$7^2$  & 5280 & 5280 & 9555 & $32^2$ & 647548785 & 5626199040 & 6457085712 \\
$8^2$  & 836559012432 & 77545495104000 & 1563557678865 & $33^2$ & 76396812236640 & 145116788951160 & 579009729388552 \\
$9^2$  & 1951960680 & 16612374240 & 39358195240 & $34^2$ & 3885 & 73920 & 73920 \\
$10^2$ & 2400 & 4389 & 8736 & $35^2$ & 121275 & 315744 & 1154400 \\
$11^2$ & 4844280 & 7134120 & 23746008 & $36^2$ & 300986505 & 660694320 & 776624128 \\
$12^2$ & 257920 & 266112 & 3932145 & $37^2$ & 110590166232 & 193219146120 & 738057963240 \\
$13^2$ & 10440 & 15960 & 62832 & $38^2$ & 804196431456 & 7816246018581 & 1829965027200 \\
$14^2$ & 292485 & 487008 & 2328480 & $39^2$ & 160888 & 1318680 & 8172360 \\
$15^2$ & 157326624 & 14657944675 & 183629376 & $40^2$ & 32367246681 & 161813106000 & 566840338944 \\
$16^2$ & 175305 & 802560 & 929040 & $41^2$ & 68986596728 & 91584309960 & 92354600520 \\
$17^2$ & 45936 & 134400 & 249711 & $42^2$ &   &   &   \\
$18^2$ & 294525 & 655776 & 2855776 & $43^2$ & 839040 & 4362072 & 4462920 \\
$19^2$ & 11898227880 & 11898227880 & 15944365080 & $44^2$ &   &   &   \\
$20^2$ & 639600 & 2246601 & 3489024 & $45^2$ & 538200 & 1065064 & 2952936 \\
$21^2$ & 17412649408 & 246163829760 & 1357746969735 & $46^2$ & 3360 & 3360 & 7293 \\
$22^2$ & 66617760 & 688642080 & 178836645 & $47^2$ & 13167 & 349440 & 526320 \\
$23^2$ & 41496 & 138600 & 138600 & $48^2$ & 13503321 & 58489600 & 98960400 \\
$24^2$ & 166512640 & 184524480 & 6527508273 & $49^2$ & 531960 & 1785168 & 3796200 \\
$25^2$ & 2429856 & 3377619 & 10594400 & $50^2$ & 39525 & 384384 & 1329216 \\
\hline
\end{tabular}
\caption{Hilbert cubes $H(a_{0}; a_{1}, a_{2}, a_{3})\subset\cal{S}$ with $a_{0}=n^2, n\in\{1,\ldots, 50\}$.}
\end{table}
\end{landscape}

\begin{landscape}
\begin{table}[htbp]
\centering
\begin{tabular}{|l|l|l|l||l|l|l|l|}
\hline
$a_{0}$ & $a_{1}$ & $a_{2}$ & $a_{3}$ &
$a_{0}$ & $a_{1}$ & $a_{2}$ & $a_{3}$ \\
\hline
$51^2$ & 166320 & 166320 & 2120248 & $76^2$ & 549114862234224 & 4122693075650625 & 20196553848570624 \\
$52^2$ & 11300325105 & 31353073920 & 33667843440 & $77^2$ & 12195120 & 74419200 & 887319015 \\
$53^2$ & 21216 & 28875 & 153216 & $78^2$ & 748672960 & 1037735712 & 10501529445 \\
$54^2$ & 9790308000 & 44245016109 & 68817025984 & $79^2$ & 286440 & 286440 & 901968 \\
$55^2$ & 15057741075 & 70139692896 & 92139564000 & $80^2$ &   &   &   \\
$56^2$ &   &   &   & $81^2$ & 6758640 & 28851823 & 76573440 \\
$57^2$ & 708142072 & 7438541760 & 12494094480 & $82^2$ & 31051605 & 182189280 & 241422720 \\
$58^2$ & 44160 & 94605 & 110880 & $83^2$ & 207480 & 337680 & 1416360 \\
$59^2$ & 25080 & 65688 & 93240 & $84^2$ & 1028093040 & 4312147833 & 12967736320 \\
$60^2$ & 255996400 & 1462599936 & 1795637025 & $85^2$ & 421800 & 4232256 & 7905744 \\
$61^2$ & 21983000 & 3335636304 & 3802075200 & $86^2$ & 27387360 & 30319653 & 54073920 \\
$62^2$ & 6765 & 43680 & 43680 & $87^2$ & 27142672471200 & 51905252513331 & 65136117775456 \\
$63^2$ & 7108920 & 16052080 & 33070032 & $88^2$ & 1414053072 & 5534161920 & 6435721985 \\
$64^2$ & 62985 & 270480 & 346368 & $89^2$ & 1899240 & 1899240 & 2299440 \\
$65^2$ & 9488336 & 19505664 & 44618175 & $90^2$ & 8029125 & 21763456 & 126824544 \\
$66^2$ & 3659040 & 4165408 & 4752405 & $91^2$ & 16728000 & 18559200 & 40568619 \\
$67^2$ & 47040 & 86112 & 126555 & $92^2$ & 110851377 & 149026800 & 637655040 \\
$68^2$ & 54933120 & 60259545 & 82514432 & $93^2$ & 200200 & 407376 & 488376 \\
$69^2$ & 51408 & 196840 & 344520 & $94^2$ & 200928 & 1551165 & 3109920 \\
$70^2$ & 36309 & 79200 & 114816 & $95^2$ & 7125883200 & 8042493375 & 213369153536 \\
$71^2$ & 1459311360 & 1730222175 & 3062842608 & $96^2$ & 32875482640 & 121549143105 & 209822618880 \\
$72^2$ & 10320319737 & 14224850640 & 178225815040 & $97^2$ & 1294755 & 2015520 & 12728352 \\
$73^2$ & 29640 & 224112 & 304920 & $98^2$ & 118560 & 908160 & 1456917 \\
$74^2$ & 405405 & 6723360 & 20082848 & $99^2$ & 45674280 & 7131286008 & 929691280 \\
$75^2$ & 1307691 & 4209184 & 7365600 & $100^2$ & 309225 & 354816 & 836400 \\
\hline
\end{tabular}
\caption{Hilbert cubes $H(a_{0}; a_{1}, a_{2}, a_{3})$ with $a_{0}=n^2, n\in\{51,\ldots, 100\}$.}
\end{table}
\end{landscape}

\begin{conj}
For each $n\in\N_{+}$, there is a Hilbert cube $H(a_{0}; a_{1}, a_{2}, a_{3})\subset\cal{S}$ of dimension 3 with $a_{0}=n^2$.     
\end{conj}

\bigskip
In Theorem \ref{thm:H3} we proved that $H_{3}(N)\gg N^{1/8}$. Here is the picture of the behavior of $H_{3}(n)/n^{1/8}$ and $H_{3}(n)/n$ for $n\leq 10^6$. From these pictures one can conjecture that the following equalities holds:
$$
\lim_{n\rightarrow +\infty}\frac{H_{3}(n)}{n^{1/8}}=+\infty,\quad \quad\lim_{n\rightarrow +\infty}\frac{H_{3}(n)}{n}=0.
$$
\begin{figure}[!ht]
%\begin{figure}[h!]
    \centering
    \includegraphics[width=1\linewidth]{Pic_4.jpg}
    \caption{Plot of the functions $H_{3}(n)/n^{1/8}$ (left) and $H_{3}(n)/n$ (right) for $n\leq 10^6$.}
    \label{fig:Pic4}
\end{figure}

Moreover, these figures suggest the following.

\begin{ques}
What is the true order of magnitude of the function $H_{3}(N)$?
\end{ques}

To get an idea of the expected behaviour, we used the {\bf FindFit} procedure in Mathematica 14.2 \cite{Wol}. The command
\begin{center}
{\bf FindFit}[data, expression, parameters, variable]
\end{center}
determines the numerical values of the parameters that make the given expression provide the best fit to the data as a function of the specified variable.
In our case, the data set consists of the values of $H_{3}(N)$ for $N \in {1, \ldots, 10^6}$, the expression is $F(x)=a x^{b}$, the parameters are $a$ and $b$, and the variable is $x$.
For this setup, the procedure finds that the best-fit parameters are $a = 0.01798344440392967$ and $b = 0.6220626625113858$. The plots of $H_{3}(n)$ and $F(n)$ and the absolute value of the difference are given below. In the considered range the function $F(n)$ approximate $H_{3}(n)$ quite well because 
$$
\max\{|H_{3}(n)-F(n)|:\;n\leq 10^7\}=15.474
$$
This experimental approach suggest that it is reasonable to expect that $H_{3}(N)\gg N^{b}$ for some $b>0.6$ and $N\gg 1$.

\begin{figure}[!ht]
%\begin{figure}[h!]
    \centering
    \includegraphics[width=0.9\linewidth]{Pic_5.jpg}
    \caption{The plot of the functions $H_{3}(n)$ (blue) and the fitted function $F(n)=0.01798n^{0.62206}$ (yellow) for $n\leq 10^7$.}
    \label{fig:Pic4}
\end{figure}

\begin{figure}[!ht]
%\begin{figure}[h!]
    \centering
    \includegraphics[width=0.9\linewidth]{Pic_6.jpg}
    \caption{The plot of $|H_{3}(n)-F(n)|$, where $F(n)$ is the fitted function $F(n)=0.01798n^{0.62206}$ for $n\leq 10^7$.}
    \label{fig:Pic4}
\end{figure}
\bigskip



Having infinitely many Hilbert cubes of dimension 3 in squares, it is natural to return to the question whether there exists a reduced Hilbert cube of dimension 4 in the set of squares (the first part in Question \ref{firstques})? Although we tried hard, we were unable to do so. However, using the parametrization given by (\ref{a1a2solution}), we were able to find the following:
\begin{align*}
    a_0 = & (c-1)^4(7c^4 + 12c^3 - 22c^2 + 12c + 7)^2, \\
a_1 = & a_2 = -24(c-1)^4(c+1)^2(c^2+1)(c^2 - 6c + 1)(c^2 + 6c + 1), & \\
a_3 = & -4c(c-1)^2(c^2-5)(c^2+7)(5c^2-1)(7c^2+1), \\
a_4 = & -\frac{1}{4}((c^3+c^2+19c-5)(3c^3-c^2+9c+5)(5c^3-19c^2-c-1)(5c^3+9c^2-c+3)
\end{align*}
where thirteen of the sixteen sums are squares (the exceptions being $a_0+a_4$, $a_0+a_1+a_2+a_4$, $a_0+a_1+a_2+a_3+a_4$).

Numerically, we can do a little better. Suppose given a Hilbert cube
$\{a_0;a_1,a_1,a_3\}$. Its extension to a $4-$cube demands finding $a_4=X$ satisfying that all the six quantities
\[ a_0+X, \; a_0+a_1+X, \; a_0+2a_1+X, \; a_0+a_3+X, \; a_0+a_1+a_3+X, \; a_0+2a_1+a_3+X, \]
be square. Consider (for example) the elliptic curve
\[ Y^2 = (a_0+X)(a_0+a_1+X)(a_0+2a_1+X), \]
with Mordell-Weil group $\mathcal{G}$. The group will have rank at least 1 since $X=0$ gives a point of infinite order. The points $(X,Y)$ with $a_0+X$, $a_0+a_1+X$, $a_0+2a_1+X$ all square are precisely the points of $2\, \mathcal{G}$. This guarantees finding $X$ with three of the above six quantities square. Let $X$ range over $X$-coordinates of $2\mathcal{G}$, and check to see whether any of
\[ a_0+a_3+X, \; a_0+a_1+a_3+X, \; a_0+2a_1+a_3+X  \]
are square. This was implemented for the parametrized $3$-cubes at (\ref{a1a2solution}) above, for small values of $c,d$. Unfortunately only a couple of examples with one extra square were found. So at present we know only ``pseudo-"$4$-cubes where fourteen of the sixteen sums are squares, e.g.:
\[ 6310^2, \; 105386400, \; 105386400, \; 144545984, \; -121397859. \]

\bigskip

Another question which comes to mind is what kind of result can be proved in the case of other polynomials of degree 2. More precisely, let $f\in\Q[x]$ be an polynomial of degree 2. Without loss of generality we can assume that for some rational numbers $a, b, a>0$, we have $f(x)=ax^2+bx$. Then, one can ask what is the biggest dimension of a Hilbert cube sitting in the set 
$$
S(a,b):=\{f(n):\;n\in\N\}\cap\N_{+}.
$$
Let us note that if $f(x)=x(x+1)/2$, i.e., $a=b=1/2$, then there are infinitely many Hilbert cubes of dimension 3 in $S(a,b)$. More precisely, for each $n\in\N$ we have $H(a_{0}(n); a_{1}(n), a_{2}(n), a_{3}(n))\subset S(1, 0)$, where
$$
a_{0}(n)=\frac{n(n+1)}{2}, \quad a_{1}(n)=66(2n+1)^2,\quad a_{2}(n)=a_{3}(n)=105(2n+1)^2.
$$
In consequence, the number of Hilbert cubes sitting in the set $S(1,0)\cap [0,N]$ for large $N$ is $\gg N^{1/2}$.

\begin{prob}
Find values $a, b$ such that the set $S(a,b)$ contains Hilbert cube of dimension 4.
\end{prob}



%The case $a_0=0$ is even for $d=3$ a famous open problem, the so called 
%cuboid problem, asking if there exists a cuboid where all side lengths, and all %planar diagonals and the spatial diagonal are all integers.

%\bigskip

%What can be proved (or what is known) about Hilbert cubes in the set


%\begin{ques}
%Let $a\in\Z\setminus\{0\}$. Does the set
%$$
%\cal{S}_{a}:=\{x_{1}^2+ax_{2}^2:\;x_{1}, x_{2}\in\N\}
%$$
%contain Hilbert cubes of arbitrarily large dimension? an infinite Hilbert cube? 

%We note that a result of the second author and Frei \cite{Elsholtz-Frei} implies %that for each binary quadratic form the value set has arithmetic progressions of %bounded length only.

%\end{ques}



%Exclude infinite Hilbert cubes in binary quadratic forms.\\


\begin{thebibliography}{99}

\bibitem{Alsetri-Shao}
A.~Alsetri, Xuancheng  Shao,
On Hilbert cubes and primitive roots in finite fields.
Arch. Math. (Basel) 118 (2022), no. 1, 49--56. 

\bibitem{AlonandAngelandBenjaminiandLubetzky:2012}
N.~Alon, O.~Angel, I.~Benjamini, E.~Lubetzky,
Sums and products along sparse graphs.
Israel Journal of Mathematics 188 (1) (2012), 353--384.

\bibitem{BombieriGranvilleandPintz}
E.~Bombieri, A.~Granville, J.~Pintz, Squares in arithmetic progressions. Duke Math. J., 66, 1992, 369--385.

\bibitem{BombieriandZannier}
E.~Bombieri, U.~Zannier, A note on squares in arithmetic progressions. II.
Atti Accad. Naz. Lincei Cl. Sci. Fis. Mat. Natur. Rend. Lincei (9) Mat. Appl. 13 (2002), no. 2, 69--75. 



\bibitem{Mag} Wieb Bosma, John Cannon, and Catherine Playoust. The Magma algebra system. I. The user language.  J. Symbolic Comput., 24(3-4) (1997), 235--265.

\bibitem{Bremner:1999}
A.~Bremner,
On squares of squares.
Acta Arith. 88 (1999), no. 3, 289--297. 

\bibitem{Bremner:2001}
A.~Bremner,
On squares of squares. II.
Acta Arith. 99 (2001), no. 3, 289--308. 

\bibitem{BrownandErdosandFreedman:1990}
T.C.~Brown, P.~Erd\H{o}s, A.R.~Freedman,
Quasi-progressions and descending waves. 
J. Comb. Theory, Ser. A 53 (1990), no.~1, 81--95.

\bibitem{CaporasoHarrisandMazur}
 L. Caporaso, J. Harris, and B. Mazur, Uniformity of rational points, J. Amer. Math. Soc., 10 (1997), pp. 1--35.
\bibitem{CillerueloandGranville:2007}
J.~Cilleruelo, A.~Granville,
Lattice points on circles, squares in arithmetic progressions, 
and sumsets of squares.
In ``Additive Combinatorics", CRM Proceedings \& Lecture Notes,
vol. 43 (2007), 241--262.

\bibitem{Conlon-Fox-Sudakov:2014}
     D.~Conlon, J.~Fox and B.~Sudakov, Short proofs of some extremal results, Combinatorics, Probability and Computing 23 (2014), 8--28.


\bibitem{Dietmann-Elsholtz:2012}
R.~Dietmann, C.~Elsholtz, 
Hilbert cubes in progression-free sets and in the set of squares.
Israel Journal of Mathematics, 
Israel J. Math. 192 (2012), 59--66.

\bibitem{Dietmann-Elsholtz:2015}
R.~Dietmann, C.~Elsholtz, 
Hilbert cubes in arithmetic sets.
Rev. Mat. Iberoam. 31 (2015), no. 4, 1477--1498. 

\bibitem{Dietmann-Elsholtz-Shparlinski}
R.~Dietmann, C.~Elsholtz, I.E. Shparlinski,
Prescribing the binary digits of squarefree numbers and quadratic residues.
Trans. Amer. Math. Soc. 369 (2017), no. 12, 8369--8388. 

\bibitem{Dujella:2001}
A.~Dujella,
On the number of Diophantine m-tuples. Ramanujan J. 15 (2008), no. 1, 37--46. 

\bibitem{DujellaandElsholtz:2013}
A.~Dujella, C.~Elsholtz, Sumsets being squares,
Acta Math. Hungar. 141 (2013), 353--357. 

\bibitem{Elsholtz-Frei}
C.~Elsholtz, C.~Frei, 
Arithmetic progressions in binary quadratic forms and norm forms.
Bull. Lond. Math. Soc. 51 (2019), no. 4, 595--602. 

\bibitem{Elsholtz-Ruzsa-Wurzinger}
C.~Elsholtz, I.Z.~Ruzsa, L.~Wurzinger
Sumset growth in progression-free sets, to appear in Acta Arith.

\bibitem{Elsholtz-Wurzinger}
C.~Elsholtz, L.~Wurzinger,
Sumsets in the set of squares.
Q. J. Math. 75 (2024), no. 4, 1243--1254. 
\bibitem{Gyarmati:2001}
K.~Gyarmati, On a problem of Diophantus.
Acta Arith. 97 (2001), no. 1, 53--65. 


\bibitem{Hegyvari:1999}
N.~Hegyv\'{a}ri, 
On the dimension of the Hilbert cubes. J. Number Theory 77 (1999), 326--330.

\bibitem{HegyvariandSarkozy:1999}
N.~Hegyv\'ari, A.~S\'ark\"ozy, On Hilbert cubes in certain sets. The
Ramanujan Journal 3 (1999), 303--314.


\bibitem{Hilbert:1892}
D.~Hilbert, \"{U}ber die Irreduzibilit\"{a}t 
ganzer rationaler Funktionen mit ganzzahligen Koeffizienten.
J. Reine Angew. Math. 110 (1892), 104--129.




\bibitem{PARI2} The PARI~Group, PARI/GP version \texttt{2.13.4}, Univ. Bordeaux, 2022, \url{http://pari.math.u-bordeaux.fr/}.

\bibitem{RivatSarkozyandStewart}
J.~Rivat, A.~S\'{a}rk\"{o}zy, C.L.~Stewart,
Congruence properties of the $\Omega$-function on sumsets.
Illinois J. Math. 43 (1999), no. 1, 1--18. 

\bibitem{Robertson}
J.P.~Robertson,
Magic squares of squares.
Math. Mag. 69 (1996), no. 4, 289--293.

\bibitem{Schoen}
T.~Schoen,
Arithmetic progressions in sums of subsets of sparse sets.
Acta Arith. 147 (2011), no. 3, 283--289. 


\bibitem{ShkredovandSolymosi}
I.~Shkredov, J.~Solymosi,
The uniformity conjecture in additive combinatorics.
SIAM J. Discrete Math. 35 (2021), no. 1, 307--321. 


\bibitem{Solymosi:2007}
J.~Solymosi, Elementary Additive Combinatorics. 
In ``Additive Combinatorics'', CRM Proceedings \& Lecture Notes,
vol.~43 (2007), 29--38.

\bibitem{Szemeredi:1969}
E.~Szemer\'{e}di,
On sets of integers containing no four elements in arithmetic progression.
Acta Math. Acad. Sci. Hung. 20 (1969), 89--104. 



\bibitem{Wol} Wolfram Research, Inc., Mathematica, Version 14.2. (Champaign, IL, 2024).



\end{thebibliography}

\end{document}



